\documentclass[11pt,reqno]{amsart}
\usepackage[margin=1.15in]{geometry}
\usepackage{amsmath,amssymb,amsthm}
\usepackage{booktabs}
\usepackage{graphicx}
\usepackage{longtable}
\usepackage{array}
\usepackage[hidelinks]{hyperref}

\theoremstyle{plain}
\newtheorem{theorem}{Theorem}
\newtheorem{proposition}[theorem]{Proposition}
\newtheorem{lemma}[theorem]{Lemma}
\newtheorem{corollary}[theorem]{Corollary}
\theoremstyle{definition}
\newtheorem{definition}[theorem]{Definition}
\newtheorem{example}[theorem]{Example}
\newtheorem{question}[theorem]{Question}
\theoremstyle{remark}
\newtheorem{remark}[theorem]{Remark}

\newcommand{\Z}{\mathbb{Z}}
\newcommand{\Q}{\mathbb{Q}}
\newcommand{\R}{\mathbb{R}}
\newcommand{\T}{\mathbb{T}}
\newcommand{\fp}[1]{\{#1\}}
\DeclareMathOperator{\Orb}{Orb}

\begin{document}

\title[On a problem of Eugene Lawler]{On a problem of Eugene Lawler:\\
knot counts in an interval deletion game}

\author{D. V. Feldman}
\address{Department of Mathematics and Statistics, University of New Hampshire,
Durham, New Hampshire 03824}
\email{dvfinnh@gmail.com}

\date{25 August 2026}

\subjclass[2020]{Primary 37B10; Secondary 11K16, 11R06, 68R05}
\keywords{Pisot number, plastic number, symbolic dynamics,
finite automaton, doubling map, Bernoulli convolution, transversality,
Mahler's problem, formal verification}

\begin{abstract}
A strip is repeatedly shortened by deleting a subinterval occupying a fixed
proportion of its length, closing the gap, and rescaling to the original
length. Deleting from either end simply truncates the strip, but an interior
deletion joins two pieces and leaves a persistent mark, which we call a knot;
later deletions may excise earlier knots. Writing $\lambda>1$ for the reciprocal
of the retained proportion, we ask how many knots can be present at once.

We show that the maximum knot count is non-decreasing in the number of moves
for every $\lambda$, and that when $\lambda$ is a Pisot number --- an
algebraic integer greater than $1$ whose other roots all have modulus less
than $1$ --- the forward orbit of the point at which knots are born is finite,
so that the game reduces
to a finite automaton and the maximum is attained and computable. We determine
that maximum for the golden ratio, where it equals $2$, and for the plastic
number, where it equals $7$. For $\lambda=\sqrt2$, $\lambda=3/2$ and
$\lambda=\sqrt3$ no such reduction is available; we report exhaustive
computations and record what is and is not settled. A second and independent
threshold occurs at $\lambda=3/2$, where the three deleted intervals tile the
strip. For $\lambda=\sqrt2$ the identity $\lambda^2=2$ makes pairs of moves act
as the two branches of the doubling map, reducing the survival of a knot to the
binary expansion of $b\sqrt2$ for an associated integer $b$. For $\lambda=3/2$ the three deleted intervals are exactly the three ternary
cells, and a knot survives precisely as long as an associated number, formed by
repeatedly multiplying by $3/2$ and subtracting an integer, remains inside a
fixed window. That is the recursion appearing in Mahler's unsolved problem on
the powers of $3/2$; the difference is that here several such numbers must be
held inside the window at once, by a single shared sequence of subtractions,
and it is this sharing rather than the width of the window that carries the
difficulty. At $\lambda=(1+\sqrt2)/2$, which has a small conjugate root but is not an
algebraic integer, the orbit of the birth point is infinite -- so the Pisot
hypothesis is sharp -- and the record lengths grow quadratically as far as they
are known. Allowing $\lambda$ to vary, we prove that two distinct knot histories, viewed
as functions of $\lambda$, always separate at a definite rate wherever their
values nearly agree; this holds on a range of $\lambda$ containing $3/2$, and it bounds on average
how many pairs of histories can nearly collide. That it could hold there was a
genuine question: the histories are encoded by power series with coefficients
$-1$, $0$, $1$, and if the coefficients are instead allowed to be arbitrary
numbers between $-1$ and $1$ the separation statement becomes false just short
of $3/2$ --- so the result rests essentially on the coefficients being
discrete. With a lower bound on the number of
survivable move sequences, and lemmas describing how those sequences branch
below the golden ratio, this supplies the first steps --- all machine-verified
--- of a programme aimed at showing the count unbounded for almost every
$\lambda$.
\par
We also show that the maximum is exactly $1$ when the
deleted interval occupies at least half the strip, and that a periodic run can never carry a knot back to
its birth point, so that unboundedness cannot be certified by an orbit that is
periodic -- though whether some periodic run keeps a knot alive forever, which
would give unboundedness at once, remains open.
\end{abstract}

\maketitle

\tableofcontents


\section{The game}\label{sec:game}

The problem treated here is recalled as having circulated among theoretical
computer scientists at Berkeley, attributed to Eugene Lawler. No published
source has been located, and what follows is a reconstruction. We take the
original to be the case in which the deleted piece is a third of the strip ---
so that the three possible deletions exactly tile it, which is what makes the
game natural to pose --- and it is this case, $\lambda=3/2$, that we regard as
the problem proper; the one-parameter family in which $\lambda$ may be any
number greater than $1$ is our own, and serves both to locate what is special
about $3/2$ and to supply cases that can be settled. A reader interested only
in the original question may read Sections~\ref{sec:game},
\ref{sec:normal}, \ref{sec:gaps} and~\ref{sec:threehalves}.

Fix $\lambda>1$ and set
\[
  r=\lambda^{-1},\qquad g=1-r .
\]
The strip is normalised at every stage to the interval $[0,1]$. A
\emph{configuration} is a finite subset of the open interval $(0,1)$, whose
elements are called \emph{knots}. Three \emph{moves} are available, each of
which deletes a closed subinterval of length $g$, closes the gap, and rescales
by $\lambda$. Two of them delete from an end and merely truncate; the third
deletes from the middle and joins two pieces:
\begin{itemize}
\item[$R$:] delete $[r,1]$. A knot at $x$ survives if and only if $x<r$, and
  moves to $\lambda x$.
\item[$L$:] delete $[0,g]$. A knot at $x$ survives if and only if $x>g$, and
  moves to $\lambda x-(\lambda-1)$.
\item[$M$:] delete $[r/2,\,1-r/2]$. A knot at $x$ survives if and only if
  $x<r/2$ or $x>1-r/2$; in the first case it moves to $\lambda x$, in the second
  to $\lambda x-(\lambda-1)$. The two exposed ends are joined, and the join is a
  new knot, which after rescaling lies at $1/2$.
\end{itemize}
The moves $L$ and $R$ remove an end of the strip and create no knot; only $M$
does. The deleted interval is closed: a knot lying on its boundary is destroyed.
This is the correct convention rather than a choice, since such a knot would
otherwise be carried to an endpoint of the shortened strip, which is a free end
and not a join of two pieces.

Throughout, $\fp{t}$ denotes the fractional part of $t$. Write
$f_0(x)=\lambda x$ and $f_1(x)=\lambda x-(\lambda-1)$ for the two branch
maps, so that every surviving knot moves by $f_0$ if it lies below the deleted
interval and by $f_1$ if it lies above.

A \emph{run} is a word in $\{L,M,R\}^{*}$, applied left to right to the empty
configuration. Let $N_\lambda(n)$ denote the largest number of knots present
after a run of length $n$, and
\[
  d_\lambda(k)=\min\{\,n: N_\lambda(n)\geq k\,\}
\]
the least length of a run producing $k$ simultaneous knots. The original
question is whether $N_\lambda$ is bounded.

In fact distinct knots never come to occupy the same point, so no convention
about coincidence is needed.

\begin{lemma}\label{lem:distinct}
Let $x<y$ be knots that both survive a move. Then their images are distinct.
Moreover no surviving knot is carried to $1/2$ by the move $M$, so the knot that
$M$ creates differs from every survivor. Consequently the number of knots after
a move is the number of survivors, together with one more when the move is $M$.
\end{lemma}

\begin{proof}
Under $L$ or $R$ both knots are moved by the same injective map, as they are
under $M$ when both lie below the deleted interval or both above it. In the
remaining case $x<r/2$ and $y>1-r/2$ hold strictly, so $y-x>1-r=g$ and
\[
  f_1(y)-f_0(x)=\lambda(y-x)-(\lambda-1)>\lambda g-(\lambda-1)=0 .
\]
Finally $f_0(x)=1/2$ forces $x=r/2$ and $f_1(x)=1/2$ forces $x=1-r/2$, and both
of these lie in the interval deleted by $M$.
\end{proof}

Under the opposite convention, in which a knot on the boundary of the deleted
interval is spared, coincidence does occur: the two boundary points are then
carried to $1/2$, where the newly created knot already lies. This is why the two
conventions agree at $\lambda=\varphi$ in Section~\ref{sec:pisot}.



\begin{figure}[ht]
\centering
\includegraphics{game.pdf}
\caption{The game at $\lambda=3/2$, where the three deleted intervals are
exactly the three ternary cells $[0,\frac13]$, $[\frac13,\frac23]$,
$[\frac23,1]$. The run $MRML$ is played from the empty configuration. The first
$M$ creates a knot at $\frac12$; $R$ carries it to $\frac34$; the second $M$
spares it, since $\frac34>\frac23$, carrying it to $\frac58$ while creating a
second knot at $\frac12$; and $L$ carries the pair to $\frac7{16}$ and
$\frac14$. Two knots in four moves --- one fewer move than the minimum for a
third, by Proposition~\ref{prop:lowerbound}. Had the move after either birth
been $M$, the newborn would have died, since $\frac12$ always lies strictly
inside the interval deleted by $M$.}
\label{fig:game}
\end{figure}

\section{Directions of attack}\label{sec:directions}

The question just posed --- is $N_\lambda$ finite? --- is not answered in this
paper for any $\lambda$ that is not a Pisot number, and in particular not for
$\lambda=3/2$, which is where the problem began. What follows is therefore not
a solution but a map of partial ones. Each buys something by relaxing the goal
in a different direction, and the directions are largely independent: a reader
may enter at any row of the table below and follow that thread alone.

\medskip
\begin{longtable}{@{}p{0.26\textwidth}p{0.45\textwidth}p{0.17\textwidth}@{}}
\toprule
\textbf{Relaxation} & \textbf{What it buys, and what it costs} &
\textbf{Where}\\
\midrule
\endfirsthead
\toprule
\textbf{Relaxation} & \textbf{What it buys, and what it costs} &
\textbf{Where}\\
\midrule
\endhead
\bottomrule
\endfoot
\bottomrule
\endlastfoot
Restrict the parameter &
Complete answers. $N_\lambda=1$ for $\lambda\geq2$; for four Pisot parameters
the maximum is computed exactly ($2$, $3$, $4$, $7$). The price is that the
Pisot numbers are countable, and the finiteness of the orbit that makes them
tractable is exactly what a general $\lambda$ lacks. &
\S\ref{sec:pisot}, \S\ref{sec:thresholds}\\[2pt]
Ask for less than a bound &
Exact least lengths $d_\lambda(k)$, and structural laws relating them, in
place of the supremum. The price is that no such statement constrains the
limit. &
\S\ref{sec:normal}, \S\ref{sec:nonpisot}\\[2pt]
Accept evidence &
Exhaustive tables, a quadratic growth law at $(1+\sqrt2)/2$ fitting every
known value, and measured densities matching the counting-optimal rate. The
price is that none of it is proof. &
\S\ref{sec:nonpisot}, App.~\ref{app:explore}\\[2pt]
Change the conclusion &
Ask for one run with infinitely many simultaneous knots rather than for
unbounded finite counts; this is equivalent to a checkable condition on
finite witnesses, and is the natural form of the backward game. The price is
that it is formally stronger, and the converse implication is open. &
\S\ref{sec:density}\\[2pt]
Assume a hypothesis &
A knot survives a stretch of play only if the moves are kind to it, and we call
a run \emph{kind} when the knot born at the outset survives every move of it
(Section~\ref{sec:kind}). Each finite kind run ends with its knot somewhere in
the strip. If those ending positions are dense, unboundedness follows; if in
addition short runs already reach everywhere, so does an exponential bound on
$d_\lambda(k)$. The price is that the
hypothesis is unproved at every non-Pisot parameter, though it is measured to
hold at the best possible rate. &
\S\ref{sec:density}\\[2pt]
Reduce to a known problem &
At $\lambda=3/2$ a single knot survives exactly as long as a number, repeatedly
multiplied by $3/2$ with an integer subtracted each time, stays inside a fixed
window --- the recursion of Mahler's unsolved problem on the powers of $3/2$.
What is new here is that several such numbers must be kept inside the window
simultaneously, by one and the same sequence of subtractions. At
$\lambda=\sqrt2$ two moves together act as doubling, and survival becomes a
condition on binary digits. The price is that the problems we arrive at are
themselves open: the difficulty is explained and relocated, not removed ---
though the relocation at $3/2$ is exact, an equivalence checked by machine. &
\S\ref{sec:threehalves}, \S\ref{sec:sqrt2}\\[2pt]
Aim at generic parameters &
Let $\lambda$ vary. The final position of a knot is then a function of
$\lambda$, and two different histories give two such functions; we show that
wherever they nearly agree they are pulling apart at a definite rate, so any
one pair can nearly collide only over a short range of $\lambda$. Summed over
all pairs this bounds, on average in $\lambda$, how much collision is
possible --- the estimate an argument about typical $\lambda$ requires. It
holds on a range containing $3/2$. The pairs of histories are encoded by power
series with coefficients $-1$, $0$, $1$; had the coefficients been allowed to
range over all of $[-1,1]$, the statement would be false just short of $3/2$,
so the discreteness of the coefficients is doing essential work. We also bound the number of survivable move sequences from below,
exponentially. The price is that the theorem about almost every $\lambda$
towards which these are steps is not proved here. &
\S\ref{sec:transversality}\\[2pt]
Reprove what is known &
A second proof of Pisot boundedness, from a minimum-separation argument rather
than an automaton: effective, uniform, and applicable where the automaton is
infeasible. The price is that its constants are far from sharp ($9$ against
$2$, $34$ against $7$). &
\S\ref{sec:gaps}\\[2pt]
Insist on certainty &
Machine verification of the great majority of what is asserted, with the
exceptions itemised. The price is only that it constrains what may be
claimed. &
App.~\ref{app:verify}\\
\end{longtable}
\medskip

Two threads cut across the table. Section~\ref{sec:conjugate} asks what governs
the \emph{size} of a bound once one exists, and finds that it is the modulus of
the conjugates rather than $\lambda$ itself. Section~\ref{sec:kind} closes one
route --- no periodic run returns a knot to its birth point --- while leaving
open the neighbouring one that would settle the problem at once, and explains
why the obvious algebraic attack on it fails.

Throughout, we try to say which of these registers a given statement belongs
to. A theorem is a theorem; a computation is labelled as one; a conditional
result carries its hypothesis in its statement; and the appendix records
exactly which claims have been checked by machine and which have not.

\section{Two normal forms}\label{sec:normal}

\subsection{Suffix decomposition}

A knot is born at $1/2$ and thereafter moves by maps depending only on its own
position. Consequently the survival of a knot born at time $t$ depends only on
the letters after $t$.

\begin{lemma}\label{lem:suffix}
For a run $w=w_1\cdots w_n$,
\[
  \#\{\text{knots after } w\}
   =\#\{\,t : w_t=M \text{ and } 1/2 \text{ survives } w_{t+1}\cdots w_n\,\}.
\]
\end{lemma}

\begin{theorem}\label{thm:mono}
$N_\lambda(n+1)\geq N_\lambda(n)$ for every $n$ and every $\lambda$.
\end{theorem}

\begin{proof}
Prepend a letter $c$ to $w$. Every index shifts by one, but every suffix is
unchanged, so by Lemma~\ref{lem:suffix} each knot of $w$ persists in $cw$; a
further knot may be gained at the front.
\end{proof}

Appending, by contrast, can cost a knot. For $\lambda=3/2$ the configuration
$\{1/4,\,7/16,\,55/64\}$ loses a knot to each of the three moves: the first
point blocks $L$, the third blocks $R$, and the second, lying in the middle
third, blocks $M$. Hence a run of minimal length attaining $k$ knots may admit
no useful extension, and Theorem~\ref{thm:mono} is recovered only by rebuilding
from the left.

The same mechanism that makes an $M$ immediately after a birth fatal gives a
universal lower bound on the record lengths.

\begin{proposition}\label{prop:lowerbound}
For every $\lambda$ and every $k$, $d_\lambda(k)\geq 2k-1$.
\end{proposition}

\begin{proof}
Since $r/2<1/2<1-r/2$, the point $1/2$ lies strictly inside the interval
deleted by $M$, so a knot born at time $t$ is destroyed if the move at $t+1$ is
$M$. If knots born at times $t_1<\dots<t_k$ are all alive at time $n$, then no
$t_{i+1}$ equals $t_i+1$, whence $n\geq t_k\geq t_1+2(k-1)\geq 2k-1$.
\end{proof}

\begin{corollary}\label{cor:recursive}
$d_\lambda(k+1)\geq d_\lambda(k)+2$.
\end{corollary}

\begin{proof}
Suppose $k+1$ knots are alive at time $n$, born at $t_1<\dots<t_{k+1}$. All of
the first $k$ are already alive at time $t_k$, so $t_k\geq d_\lambda(k)$; and
$t_{k+1}\geq t_k+2$ as in the proposition.
\end{proof}

\subsection{The itinerary of a knot}

A knot born at $1/2$ moves by $x_j=\lambda x_{j-1}-\varepsilon_j(\lambda-1)$,
where $\varepsilon_j\in\{0,1\}$ records the branch taken at its $j$th move. Call
$(\varepsilon_j)$ its \emph{itinerary}.

\begin{proposition}\label{prop:itinerary}
For every $n$,
\[
  x_n=(\lambda-1)\sum_{k\geq1}\varepsilon_{n+k}\lambda^{-k},
\]
and in particular the itinerary of a knot is exactly a $\{0,1\}$-representation
\[
  \sum_{j\geq1}\varepsilon_j\lambda^{-j}=\frac1{2(\lambda-1)} .
\]
Conversely every $\{0,1\}$-sequence satisfying this is the itinerary of a knot.
\end{proposition}

\begin{proof}
Inverting the recursion gives $x_{j-1}=r x_j+\varepsilon_j(1-r)$ with
$r=\lambda^{-1}$, so $x_n=(1-r)\sum_{k=1}^{N}\varepsilon_{n+k}r^{\,k-1}+r^Nx_{n+N}$,
and $r^N\to0$ with $x_{n+N}\in[0,1]$. Taking $n=0$ and $x_0=1/2$ and multiplying
by $r$ gives the second display, since $r/(1-r)=(\lambda-1)^{-1}$. For the
converse, the displayed tails lie in $[0,1]$ for every sequence, so they define
an admissible orbit.
\end{proof}

The number being represented is $1$ at $\lambda=3/2$, $(1+\sqrt2)/2$ at
$\lambda=\sqrt2$ and $\varphi/2$ at $\lambda=\varphi$.

\subsection{Circle normal form}

Identify the endpoints of $[0,1]$ and call the resulting point the \emph{seam};
the strip is then a circle carrying a marked point together with its knots. Each
of the three moves deletes an arc of length $g$, rejoins, rescales by $\lambda$,
and deposits a point at the join. The deposited point is the new seam when the
deleted arc contains the old seam, which is the case for $L$ and for $R$, and is
a knot otherwise, which is the case for $M$. Counting the seam as a marked point
therefore makes the rule uniform. Let
\begin{quote}
$D$: delete the closed arc of length $g$ centred at $0$, rejoin, rescale by
$\lambda$, place the join at $0$, and mark it,
\end{quote}
and let $\rho_\theta$ denote rotation of $\T=\R/\Z$ by $\theta$.

\begin{proposition}\label{prop:circle}
With the seam held at $0$,
\[
  L=D\circ\rho_{-g/2},\qquad R=D\circ\rho_{g/2},\qquad
  M=\rho_{-1/2}\circ D\circ\rho_{-1/2},
\]
and the number of knots is one less than the number of marked points.
\end{proposition}

Absorbing each post-rotation into the pre-rotation of the following move puts a
run in the form $D\rho_{\theta_n}D\rho_{\theta_{n-1}}\cdots D\rho_{\theta_1}$
with $\theta_j$ drawn from the six values $\{0,\pm g/2,1/2,1/2\pm g/2\}$. On the
Fourier side the rotations are diagonal and $D$ is a single fixed kernel, so the
choice of move enters only through a phase. At $\lambda=3/2$ the three
pre-rotations are $g/2+\{0,1/3,2/3\}$, and the phase is a character of $\Z/3$.

\begin{remark}\label{rem:jsr}
The homogeneous parts of the three move operators are positive and do not
increase total mass, so no product of them can expand --- their joint spectral
radius is at most $1$ --- and it
equals $1$, since $R$ preserves $(0,r)$, $L$ preserves $(g,1)$, and these cover
$(0,1)$ whenever $g<r$, that is whenever $\lambda<2$. A single knot can
therefore be maintained indefinitely, and any bound obtained from a linear
functional of the configuration degenerates to $N_\lambda(n)\leq n$. A useful
bound must be nonlinear in the configuration.
\end{remark}

\section{Pisot parameters}\label{sec:pisot}

\emph{Relaxation in force: the parameter is restricted. Within that
restriction the results here are complete.}

Recall that a \emph{Pisot number} is a real algebraic integer greater than $1$
all of whose other roots --- its conjugates --- lie strictly inside the unit
circle. The golden ratio is one, its conjugate being $-0.618$; so is the
plastic number, the real root of $x^3=x+1$, whose two conjugates are complex
of modulus $0.869$. The property we use is elementary and entirely
consequential: because the conjugates shrink under powering while the integer
lattice is discrete, $\lambda^n$ is very close to a whole number for large $n$,
and sums built from powers of $\lambda$ cannot spread out. Everything in this
section follows from that.

Knots born at $1/2$ take values in $\tfrac12\Z[\lambda]$, since $f_0$ and $f_1$
have coefficients in $\Z[\lambda]$. Let $\Orb(\lambda)$ denote the set of points
reachable from $1/2$ by successive applications of $f_0$ and $f_1$ along
surviving branches.

\begin{theorem}\label{thm:pisot}
If $\lambda$ is a Pisot number, then $\Orb(\lambda)$ is finite.
\end{theorem}

\begin{proof}
Let $\lambda$ have degree $d$ and conjugates $\lambda^{(2)},\dots,\lambda^{(d)}$,
all of modulus less than $1$. Applying a field embedding to $f_0$ and $f_1$
gives the maps $y\mapsto\lambda^{(i)}y$ and
$y\mapsto\lambda^{(i)}y-(\lambda^{(i)}-1)$, affine contractions with fixed
points $0$ and $1$ respectively. Any orbit of the semigroup they generate is
bounded, uniformly in the branch sequence. The map sending $x$ to the vector of
its conjugates embeds $\tfrac12\Z[\lambda]$ as a lattice, and every knot
satisfies $0<x<1$ while its conjugates remain in a fixed bounded set. A bounded
region contains finitely many lattice points.
\end{proof}

\begin{corollary}\label{cor:decide}
For Pisot $\lambda$ the reachable configurations form a finite set,
$N_\lambda$ is eventually constant, and $\sup_n N_\lambda(n)$ is computable by a
finite search.
\end{corollary}

Two further Pisot parameters round out the picture, both computed exactly in
$\Z[\lambda]$ (Appendix~\ref{app:verify}). For the tribonacci constant, the
real root of $x^3=x^2+x+1$ ($\lambda\approx1.83929$), the orbit of $1/2$ has
$7$ points, only $20$ configurations are reachable, and the maximum is exactly
$3$, attained first at length $7$ ($d(k)=1,3,7$). For the supergolden ratio,
the real root of $x^3=x^2+1$ ($\lambda\approx1.46557$), the orbit has $43$
points, $412$ configurations are reachable, and the maximum is exactly $4$
($d(k)=1,3,5,11$). Both closures are small enough that full kernel
certification appears feasible, unlike the plastic configuration closure. Note the tribonacci
maximum $3$ exceeds the golden-ratio maximum $2$ at a larger $\lambda$ --
non-monotonicity again.

\begin{remark}\label{rem:sharp}
The hypothesis of Theorem~\ref{thm:pisot} cannot be weakened to the contraction
of the conjugates alone. The parameter $\lambda=(1+\sqrt2)/2$, with minimal
polynomial $4x^2-4x-1$, has its conjugate $(1-\sqrt2)/2$ of modulus below $1$
but is not an algebraic integer: $\lambda^2=\lambda+\tfrac14$, so
$\Z[\lambda]$ contains $\Z[\tfrac12]$ and is dense, and the lattice step of the
proof fails. The orbit of $1/2$ is indeed infinite there; the numbers of
distinct points at successive depths run
$2,4,8,14,26,46,80,136,226,378,\dots$, growing by a factor of about $1.66$,
with $165{,}094$ points at depth $22$. Section~\ref{sec:nonpisot} examines this
parameter further.
\end{remark}

\subsection{The golden ratio}

\begin{theorem}\label{thm:phi}
For $\lambda=\varphi=(1+\sqrt5)/2$ one has $\sup_n N_\varphi(n)=2$, attained by
the run $MRM$.
\end{theorem}

\begin{proof}
By Theorem~\ref{thm:pisot} the orbit is finite; computation gives the five
points
\[
  p_1=1-\tfrac{\varphi}{2},\quad
  p_2=\tfrac{\varphi-1}{2},\quad
  p_3=\tfrac12,\quad
  p_4=\tfrac{3-\varphi}{2},\quad
  p_5=\tfrac{\varphi}{2},
\]
with approximate values $0.190983$, $0.309017$, $0.5$, $0.690983$, $0.809017$.
Here $p_2=r/2$ and $p_4=1-r/2$ are exactly the endpoints of the interval deleted
by $M$. Each transition below is an instance of $\varphi^2=\varphi+1$; for
example $\varphi\,p_2=(\varphi^2-\varphi)/2=1/2=p_3$.

\begin{center}
\begin{tabular}{lccccc}
\toprule
 & $p_1$ & $p_2$ & $p_3$ & $p_4$ & $p_5$\\
\midrule
$R$ & $p_2$ & $p_3$ & $p_5$ & --- & ---\\
$L$ & --- & --- & $p_1$ & $p_3$ & $p_4$\\
$M$ & $p_2$ & --- & --- & --- & $p_4$\\
\bottomrule
\end{tabular}
\end{center}

A dash records destruction. Under $M$ the points $p_2$ and $p_4$ lie on the
boundary of the deleted interval and $p_3$ lies in its interior; the move also
adjoins $p_3$. Enumerating the configurations reachable from $\emptyset$ gives
twelve, namely $\emptyset$, the five singletons, and
\[
  \{p_3,p_4\},\ \{p_2,p_3\},\ \{p_1,p_3\},\ \{p_3,p_5\},\ \{p_2,p_5\},\
  \{p_1,p_4\}.
\]
All have at most two elements, and $\{p_3,p_4\}$ is reached by $MRM$.
\end{proof}

The bound is insensitive to the treatment of the boundary. If a knot on the
boundary of the deleted interval were spared, then $p_2$ and $p_4$ would both be
carried by $M$ to $p_3$, where the newly created knot already lies, and the
twelve configurations would be unchanged.

Boundedness here is not a matter of knots dying. Under the periodic run
$(RLL)^\infty$ a knot born at $1/2$ returns to $1/2$ after every three moves and
survives for all time. The bound arises instead because $M$ is at once the only
creator of knots and a destroyer of them.

\subsection{The plastic number}

Let $\rho$ denote the real root of $x^3=x+1$, so $\rho\approx1.324718$. Its
conjugates have modulus $\rho^{-1/2}\approx0.8688$, so $\rho$ is Pisot and
Theorem~\ref{thm:pisot} applies.

\begin{proposition}\label{prop:plastic}
For $\lambda=\rho$ the orbit $\Orb(\rho)$ has $153$ points, there are $25{,}525$
reachable configurations, and $\sup_n N_\rho(n)=7$. The least lengths attaining
$k$ knots for $k\leq6$ are $1,3,5,7,12,17$.
\end{proposition}

This is a finite verification: the orbit is enumerated to closure, the induced
automaton on subsets is enumerated to closure, and the maximum is read off.

\section{Two independent thresholds}\label{sec:thresholds}

The three deleted intervals are $[0,g]$, $[\tfrac12-\tfrac g2,\tfrac12+\tfrac
g2]$ and $[1-g,1]$, all of length $g=1-\lambda^{-1}$. Their union is $[0,1]$ if
and only if $g\geq1/3$, that is if and only if $\lambda\geq3/2$. Three regimes
result.

\begin{center}
\begin{tabular}{lccl}
\toprule
$\lambda$ & value & $g$ & geometry\\
\midrule
$(1+\sqrt2)/2$ & $1.20711$ & $0.17157$ & safe set $(0.1716,0.4142)$ and its mirror\\
plastic $\rho$ & $1.32472$ & $0.24512$ & safe set $(0.2451,0.3774)$ and its mirror\\
$\sqrt2$ & $1.41421$ & $0.29289$ & safe set $(0.2929,0.3536)$ and its mirror\\
$3/2$ & $1.50000$ & $0.33333$ & the three deleted intervals tile $[0,1]$\\
$\varphi$ & $1.61803$ & $0.38197$ & overlap of length $0.0729$\\
$\sqrt3$ & $1.73205$ & $0.42265$ & overlap of length $0.1340$\\
\bottomrule
\end{tabular}
\end{center}

For $\lambda<3/2$ the \emph{safe set}
$(g,\tfrac12-\tfrac g2)\cup(\tfrac12+\tfrac g2,1-g)$ is non-empty, and a knot
lying in it survives every move. At $\lambda=3/2$ the three deleted intervals
are the three thirds of $[0,1]$, so that a move preserves every knot if and only
if one of the thirds contains no knot; by the pigeonhole principle two knots can
always be preserved, and three is the first count at which a configuration can
be blocked. For $\lambda>3/2$ the deleted intervals overlap and some points are
destroyed by two of the three moves.

This threshold is independent of the arithmetic of $\lambda$. The golden ratio
and the plastic number are both Pisot and both give bounded games, but they lie
on opposite sides of $3/2$ and their maxima are $2$ and $7$. The Pisot condition is what we can prove to
force boundedness; whether it is also necessary is open
(Section~\ref{sec:questions}). What governs the \emph{size} of the bound is
neither the Pisot condition nor the position relative to $3/2$, but a third
quantity, examined in Section~\ref{sec:conjugate}.

At the far end of the range the classification is complete. Observe that
$f_0(x)=\lambda x$ multiplies $x$ by $\lambda$, while
$1-f_1(x)=1-\lambda x+(\lambda-1)=\lambda(1-x)$ multiplies $1-x$ by $\lambda$.
Every move therefore multiplies one of $x$ and $1-x$ by $\lambda$, and both must
remain in $(0,1)$; a knot survives a move only if $\min(x,1-x)<r$.

\begin{proposition}\label{prop:big}
$N_\lambda(n)=1$ for every $n\geq1$ if and only if $\lambda\geq2$, that is, if
and only if the deleted interval occupies at least half the strip.
\end{proposition}

\begin{proof}
A knot is born at $x=1/2$, where $\min(x,1-x)=1/2$. If $\lambda\geq2$ then
$r\leq1/2$ and the condition $\min(x,1-x)<r$ fails, so the knot is destroyed by
whichever move follows its birth and no two knots coexist. If $\lambda<2$ then
$r>1/2$ and the run $MRM$ produces two knots.
\end{proof}

The threshold at $\lambda=2$ does not extend downwards into a monotone
principle. Exhaustive search to depth $16$--$18$ gives the following, where
entries other than those for $\varphi$ are lower bounds.

\begin{center}
\begin{tabular}{lcccccccc}
\toprule
$\lambda$ & $5/4$ & $8/5$ & $\varphi$ & $5/3$ & $7/4$ & $9/5$ & $19/10$ & $2$\\
$g$ & $.200$ & $.375$ & $.382$ & $.400$ & $.429$ & $.444$ & $.474$ & $.500$\\
\midrule
knots & $7$ & $4$ & $2$ & $4$ & $3$ & $4$ & $4$ & $1$\\
\bottomrule
\end{tabular}
\end{center}

Both neighbours of $\varphi$ listed here reach four knots, while $\varphi$
reaches two and stops. Outside the regime $g\geq1/2$, where the count collapses
for reasons independent of arithmetic, the classification is not determined by
$g$.

\section{Parameters that are not Pisot}\label{sec:nonpisot}

\emph{Relaxation in force: finitely many knots at a time, and computation in
place of proof. Nothing in this section bounds $N_\lambda$.}

\begin{lemma}\label{lem:overlap}
Neither $1/\sqrt2$ nor $2/3$ is a root of a non-zero polynomial with
coefficients in $\{0,\pm1\}$. By contrast $1/\varphi$ is a root of $x^2+x-1$ and
$1/\rho$ is a root of $x^3+x^2-1$.
\end{lemma}

\begin{proof}
Suppose $\sum_k a_k 2^{-k/2}=0$ with $a_k\in\{0,\pm1\}$, finitely many non-zero.
Separating even and odd indices gives $A+B/\sqrt2=0$ with $A,B$ rational, whence
$A=B=0$. Each of $A$ and $B$ is a finite $\{0,\pm1\}$-combination of distinct
powers of $\tfrac12$; if the least index with non-zero coefficient is $j$, that
term has modulus $2^{-j}$ while the remaining terms sum to less than $2^{-j}$,
so all coefficients vanish.

Suppose $\sum_k a_k(2/3)^k=0$ with $a_k\in\{0,\pm1\}$ and $a_k=0$ for $k>n$.
Multiplying by $3^n$ gives $\sum_k a_k2^k3^{n-k}=0$. Reduction modulo $2$ leaves
$a_0 3^n$, so $a_0$ is even and hence zero; dividing by $2$ and repeating gives
$a_k=0$ for all $k$.
\end{proof}

\begin{remark}\label{rem:overlap}
Lemma~\ref{lem:overlap} is the exact-overlap criterion for the iterated function
system $\{rx,\ rx+(1-r)\}$, which is the backward form of the game: a run of
length $n$ read from the right transforms the set of starting points surviving
the current suffix by these two contractions. Two compositions of equal length
coincide precisely when $r$ satisfies such a relation. Running this system with the two
contractions chosen by independent fair coin tosses gives a random point of
$(0,1)$, and the distribution of that point --- the \emph{Bernoulli
convolution} with parameter $r$ --- is the invariant measure of the system.
Whether it is spread smoothly or concentrated on a small set is the question
studied for eighty years under that name, and it turns on the same
coincidences: its Fourier transform fails to tend to zero exactly when
$\lambda$ is Pisot, by Erd\H{o}s
\cite{Erdos1939} and Salem \cite{Salem1944}; and the corresponding statement in real
space is the separation lemma of Garsia \cite{Garsia1962}. Under stronger
hypotheses than we invoke here, results of Hochman \cite{Hochman2014} show that
a dimension drop for such self-similar measures forces super-exponentially
close pairs of compositions -- so that, in the algebraic setting, exact
overlaps are the source of the drop. We do not use these results; we record
only that the game lives in the same arithmetic-overlap framework, and that
Lemma~\ref{lem:overlap} is exactly the exact-overlap condition for it. The finiteness of $\Orb(\lambda)$ in
Theorem~\ref{thm:pisot} is the extreme case of this collapse. For $\sqrt2$ and
$3/2$ Lemma~\ref{lem:overlap} shows that no collapse occurs, in agreement with
the exponential growth of the number of reachable configurations, measured at
about $1.41^n$ and $1.335^n$ respectively.
\end{remark}

\begin{figure}[ht]
\centering
\includegraphics{reachable.pdf}
\caption{The positions a knot born at $1/2$ can occupy after twenty moves, using
only $L$ and $R$; the count is shown at the right. By Lemma~\ref{lem:overlap}
distinct surviving words give distinct positions at $3/2$ and $\sqrt2$, so
these counts are exactly the numbers of surviving words, out of $2^{20}$; they
grow at the rate $2/\lambda$. At the golden ratio the orbit is the five-point
set of Theorem~\ref{thm:phi} and the reachable set cycles with period three
through $\{1/2\}$, $\{p_1,p_5\}$, $\{p_2,p_4\}$: exponentially many surviving
words, at most two places to be.}
\label{fig:reachable}
\end{figure}

The following table records exhaustive computation. An entry $\geq m$ means that
the search was complete to depth $m-1$ without attaining the count.

\begin{center}
\begin{tabular}{lccccccccccccc}
\toprule
$k$ & 1 & 2 & 3 & 4 & 5 & 6 & 7 & 8 & 9 & 10 & 11 & 12 & 13\\
\midrule
$d_{(1+\sqrt2)/2}(k)$ & 1 & 3 & 5 & 7 & 9 & 11 & 13 & 17 & 21 & 26 & 31 & 33--37 & 35--49\\
$d_{\sqrt2}(k)$ & 1 & 3 & 5 & 9 & 13 & 17 & 27 & 31 & $\geq39$ & & & &\\
$d_{3/2}(k)$ & 1 & 3 & 5 & 9 & 19 & 23 & 52 & $\geq57$ & & & & &\\[-1pt]
\multicolumn{14}{l}{\footnotesize (the configurations attaining these are drawn in Figure~\ref{fig:records32})}\\
$d_{\sqrt3}(k)$ & 1 & 3 & 10 & 29 & $\geq31$ & & & & & & & &\\
$d_{\rho}(k)$ & 1 & 3 & 5 & 7 & 12 & 17 & finite & --- & & & & &\\
$d_{\varphi}(k)$ & 1 & 3 & --- & & & & & & & & & &\\
\bottomrule
\end{tabular}
\end{center}

For $\lambda=3/2$ the entry $43$--$52$ records that no run of length at most
$42$ attains seven knots and that an explicit run of length $52$ does. For
$\lambda=\rho$ the seventh knot is attained, by
Proposition~\ref{prop:plastic}, but not within the depth searched exhaustively.

Two features of this table deserve comment. The maximum for the plastic number
is $7$, yet its least lengths for $k\leq6$ are smaller throughout than those for
$\sqrt2$: early growth carries no information about boundedness. And
$d_{3/2}$ is constant at $4$ from $n=9$ to $n=18$ and at $6$ from $n=23$ to at
least $n=43$, so a plateau is not evidence of a ceiling.

\begin{proposition}\label{prop:periodic}
Over all runs of period at most $6$, iterated for $240$ moves, the largest
number of simultaneous knots is $3$ for $\lambda=3/2$, $4$ for $\lambda=\sqrt2$
and $2$ for $\lambda=\varphi$.
\end{proposition}

Under a run of period $p$ the number of knots present at time $n$ equals, up to
edge effects, the sum over the $p$ phases of the lifetime of a knot born at that
phase, divided by $p$. It is therefore bounded if and only if no knot survives
for ever, that is if and only if $1/2$ lies outside the survivor set of the
periodic itinerary for each phase. This is a Diophantine condition on $\lambda$
rather than a formal one; at $\lambda=3/2$ it is of the type considered by
Mahler \cite{Mahler1968}, as Section~\ref{sec:threehalves} makes precise. Proposition~\ref{prop:periodic} shows only that no
short period succeeds.

\subsection{\texorpdfstring{The parameter $(1+\sqrt2)/2$}{The parameter (1+sqrt2)/2}}
\label{sub:silver}

Half the silver ratio deserves separate attention, both as the witness of
Remark~\ref{rem:sharp} and because it carries the strongest evidence of
unboundedness at any parameter examined. It shares its arithmetic shape with
$3/2$: each is $\theta/2$ for a Pisot number $\theta$, namely $3$ and
$1+\sqrt2$, each has a coefficient ring containing $\Z[\tfrac12]$, and the
number represented by an itinerary under Proposition~\ref{prop:itinerary},
$1/(2(\lambda-1))$, is $1$ in the first case and the silver ratio $1+\sqrt2$
itself in the second.

Its gap $g=3-2\sqrt2\approx0.1716$ leaves the largest safe set of the
parameters considered, of measure $6\sqrt2-8\approx0.485$. The first eleven
record lengths in the table are exact, and they satisfy
\[
  d(k)=\max\bigl(2k-1,\ \lfloor k^2/4\rfloor+1\bigr)
  \qquad(1\leq k\leq 11),
\]
the first branch, which is Proposition~\ref{prop:lowerbound}, holding with
equality up to $k=7$ and the second from $k=7$ on. Whether the formula persists
is open; it predicts $d(12)=37$ and $d(13)=43$, and what is known is
$d(12)\in[33,37]$ and $d(13)\in[35,49]$: the lower bounds from
Corollary~\ref{cor:recursive}, the upper bound $37$ from a beam search, and the
upper bound $49$ from an explicit $49$-move run with thirteen simultaneous
knots, found by a search that minimises the knots' ages rather than the run
length and checked three independent ways. If it does persist, then $N_\lambda(n)$ grows like $2\sqrt n$: polynomial
growth, against the doubling pattern of the $3/2$ column.

The remaining computations of this paper, run at this parameter, give: over all
periodic runs of period at most $6$, iterated for $240$ moves, seven
simultaneous knots, attained by $(LM)^\infty$, against three at $3/2$ and four
at $\sqrt2$; over all periodic sequences of period at most $14$ containing an
$M$, a longest survival of $1/2$ of $67$ moves, at period $11$, with none of the
$381{,}952$ surviving period-$14$ prefixes reaching $200$; a kind set of
dimension approximately $0.834$, against $0.631$, $0.679$ and $0.421$ at $3/2$,
$\sqrt2$ and $\varphi$ (Proposition~\ref{prop:kinddim}); and no Littlewood polynomial of degree at most $12$
vanishing at $\lambda$, nor any $\{0,\pm1\}$ polynomial of degree at most $10$
vanishing at $1/\lambda$, so that the conclusions of Section~\ref{sec:kind}
apply here unchanged: $1/2$ is never a periodic point, separation is
exponential, and a single periodic kind sequence would give
$N_\lambda=\infty$.

\section{Transversality and branching: certified steps toward almost every
\texorpdfstring{$\lambda$}{lambda}}\label{sec:transversality}

\emph{Relaxation in force: the target is generic rather than specific
parameters. The theorems here are steps of a programme whose conclusion is not
reached.}

The technical term for what this section proves is \emph{transversality}, and
it is worth saying in ordinary words what it asserts before the notation
begins. Fix two different histories a knot might have --- two sequences of
survived moves --- and let $\lambda$ vary. Each history determines where the
knot ends up, so each gives a function of $\lambda$, and the two knots collide
at those $\lambda$ where the two functions agree. Transversality is the
statement that this can only happen cleanly: wherever the two values are close,
the two functions are pulling apart at a definite rate, so the set of $\lambda$
for which they are within $\rho$ of each other is a single short interval, of
length proportional to $\rho$. Since there are exponentially many pairs of
histories and each contributes only a short interval, one can add up the
contributions and bound, on average over $\lambda$, how much collision occurs
in total. That average bound is what an argument about almost every $\lambda$
consumes, and it is Proposition~\ref{prop:paircount} below.

Two facts make this delicate. The first is that transversality is not
automatic: it genuinely fails for large $\lambda$. The technique available
since Solomyak's theorem certifies a wider class of series --- coefficients
anywhere in $[-1,1]$ --- and that wider class really does fail at
$\lambda\approx1.5405$, just short of $3/2$; only the smaller class our
problem actually produces, with coefficients $-1$, $0$, $1$, can go further.
Proposition~\ref{prop:restricted} is the work of showing that it does. The second
is that the pairs one needs to count are not arbitrary: they must be histories
in which the knot actually survives, and controlling those is the subject of
the second half of the section.

The density criterion of Section~\ref{sec:density} reduces unboundedness to a
property of one orbit, and Question~\ref{q:ae} asks for it almost everywhere.
This section records how far that programme has been carried, with proofs
where they are short and machine-verified certificates where they are not.
Everything stated as a proposition here is certified in Lean
(Appendix~\ref{app:verify}).

\subsection{The obstruction just short of $3/2$, and the way past it}
For branch words $\varepsilon\neq\varepsilon'$ of length $m$ first differing at
index $k$, the endpoint functions of Section~\ref{sec:density} satisfy
\[
\bigl|\Phi_\varepsilon(\lambda)-\Phi_{\varepsilon'}(\lambda)\bigr|
=(\lambda-1)\lambda^{m-k}\bigl|g(1/\lambda)\bigr|,\qquad
g(x)=1+\textstyle\sum_{i\geq1}c_ix^i,\ c_i\in\{-1,0,1\},
\]
so everything reduces to a statement about the class $\mathcal B_{01}$ of such
power series --- those with constant term $1$ and every other coefficient
$-1$, $0$ or $1$. The statement wanted is \emph{transversality}: that
$|g(x)|<\delta$ forces $g'(x)<-\delta$, so that a member of the class which is
near zero is on its way down through zero and cannot linger.

The way such statements have been proved since Solomyak's theorem
\cite{Solomyak1995} is to exhibit one explicit worst-case series and check it
alone. One takes
$h(x)=1-x-\dots-x^{k-1}+a\,x^k+x^{k+1}+x^{k+2}+\cdots$, the series that falls
as fast as the coefficients allow and then rises again; if $h$ is still
positive and still decreasing at $x_0$, then transversality holds throughout
$[0,x_0]$, because any series of the class which is small at a point is
dominated there by a translate of $h$. These are the $(*)$-functions of
Peres and Solomyak \cite{PeresSolomyak1996}.

Two things limit the argument. It certifies not our class but the larger one
in which the coefficients may be any real numbers in $[-1,1]$, and its reach
is bounded by how that larger class behaves: choosing $k=4$ and $a=0.1$ gives
$h(0.63)=0.0070$ and $h'(0.63)=-0.497$, so transversality holds on
$[0,0.63]$, and no choice pushes past $x\approx0.6487$. Nor could any other
argument, because at $x^*=0.649138\ldots$ the larger class genuinely fails:
some member of it has a double zero there, sitting on the axis instead of
crossing it. We locate $x^*$ by linear programming. Since $x=1/\lambda$, that
ceiling is $\lambda\approx1.5405$ --- and $\lambda=3/2$ sits at $x=2/3$,
outside, by about $0.017$.

The way through is that our class is smaller: coefficients $-1$, $0$, $1$ and
nothing between. That turns out to be exactly enough.

\begin{proposition}[machine-verified]\label{prop:restricted}
For every $c:\mathbb N\to\{-1,0,1\}$ with $c_0=1$ and every
$x\in[\tfrac12,\tfrac{667}{1000}]$: if $|g(x)|\leq\tfrac1{1000}$ then
$g'(x)<-\tfrac1{1000}$.
\end{proposition}

The certificate is a finite case check. The window is cut into $27$ pieces;
on each piece one bounds the possible values of $g$ and $g'$ over all
coefficient choices at once, using that a partial sum plus a bound on
everything that can follow it already excludes most cases; a case not excluded
is split further. Twenty-seven pieces suffice, and the whole check is carried
out in integer arithmetic by the proof kernel, with the pieces themselves
proved to cover the window; the class is the full infinite one, the
value a convergent series, and the conclusion holds for the honest derivative.
The window cannot extend much further, and this is now a theorem rather than a
search report. At $x=\tfrac{3339}{5000}=0.6678$, only $0.0008$ beyond the
certified window, an explicit member of the class of degree $22$ takes a value
in $(0,\tfrac1{1000}]$ while its derivative stays above $-\tfrac1{1000}$: small
without falling fast, which is exactly what the proposition forbids inside the
window. At $x=\tfrac{3343}{5000}$ another explicit member, of degree $26$, is
smaller than $10^{-5}$ and \emph{rising}. Both are kernel checks over exact
rationals (the programs that found them are not trusted; the kernel re-verifies
what they propose). So for this class and this $\delta$ the last possible
endpoint of the window is pinned between $0.667$ and $0.6678$, a band of width
less than a thousandth --- and $3/2$, at $x=2/3\approx0.6667$, clears the
demonstrated failures by about a thousandth.

\begin{proposition}[pair counting; machine-verified]\label{prop:paircount}
Let $I=[\tfrac{1000}{667},2]$, $\lambda_0=\tfrac{1000}{667}$,
$\delta=\tfrac1{1000}$. For distinct branch words of length $m$ first
differing at $k$, and every $\rho\geq0$,
\begin{align*}
\mathrm{Leb}\bigl\{\lambda\in I:|\Phi_\varepsilon-\Phi_{\varepsilon'}|\leq\rho\bigr\}
&\leq\frac{8\,\rho\,\lambda_0^{k-m}}{\delta(\lambda_0-1)},\\
\int_I\#\{\rho\text{-close pairs}\}\,d\lambda
&\leq\frac{4\lambda_0\,\rho}{\delta(\lambda_0-1)(2-\lambda_0)}
\Bigl(\frac4{\lambda_0}\Bigr)^{m}.
\end{align*}
\end{proposition}

This is the estimate a parameter-exclusion argument consumes: the parameter
mass on which any two endpoints collide at scale $\rho$ is linear in $\rho$,
uniformly over the exponentially many pairs, on a window containing $3/2$.

\subsection{Branching below the golden ratio}
The other half of the programme concerns the kind words themselves. A branch
survives from $x$ exactly when its image stays in $(0,1)$: branch $0$ needs
$x<r$, branch $1$ needs $x>g$. On the \emph{window} $(g,r)$ both branches are
legal; elsewhere exactly one is, and the dynamics is forced.

\begin{lemma}[no jumping]\label{lem:nojump}
For $\lambda^2<\lambda+1$ the forced dynamics cannot cross the window: from
$x\leq g$, either $\lambda x\leq g$ or $\lambda x\in(g,r)$; from $x\geq r$,
either $\lambda x-(\lambda-1)\geq r$ or it lies in $(g,r)$. Each crossing
inequality is equivalent to $\lambda^2\geq\lambda+1$, and for
$\lambda\geq\varphi$ the pair $\{1/(\lambda+1),\lambda/(\lambda+1)\}$ is a
forced two-cycle avoiding the window forever.
\end{lemma}

\begin{proof}
A jump from below requires $\lambda g\geq r$ and an undershoot from above
requires $\lambda r-(\lambda-1)\leq g$; each rearranges to
$\lambda^2-\lambda-1\geq0$. For the cycle, $\lambda\cdot\frac1{\lambda+1}
=\frac{\lambda}{\lambda+1}$ and $\lambda\cdot\frac{\lambda}{\lambda+1}
-(\lambda-1)=\frac1{\lambda+1}$, and $\frac1{\lambda+1}\leq g$,
$\frac{\lambda}{\lambda+1}\geq r$ are again both equivalent to
$\lambda^2\geq\lambda+1$.
\end{proof}

The golden ratio thus appears in this paper for a second time, for what looks
like a second reason: Section~\ref{sec:pisot} used $\varphi^2=\varphi+1$ to
close an orbit, and here the same equation is the exact threshold at which the
forced dynamics can leap the window. We do not know a conceptual reason the
two thresholds coincide.

\begin{lemma}[good child]\label{lem:goodchild}
Let $[\lambda_0,\lambda_1]\subset(1,\varphi)$ and
$\eta=\min(\lambda_0-1,(2-\lambda_1)/2)$. Since $g+r=1$, the midpoint of the
window is $\tfrac12$. For $x$ in the window, the child $\lambda x$ (if
$x\leq\tfrac12$) or $\lambda x-(\lambda-1)$ (if $x>\tfrac12$) lies in
$[\eta,1-\eta]$, and its forced orbit re-enters the window within
$B=\lceil\log(g/\eta)/\log\lambda_0\rceil+1$ steps.
\end{lemma}

\begin{proof}
For $x\leq\tfrac12$ the child lies in $(\lambda-1,\lambda/2]$; for
$x>\tfrac12$, in $[(2-\lambda)/2,2-\lambda)$, and $2-\lambda<r$ because
$(\lambda-1)^2>0$. In both cases the distance to $\{0,1\}$ is at least
$\eta$; a forced orbit multiplies that distance by $\lambda$ at each step and,
by Lemma~\ref{lem:nojump}, enters the window rather than crossing it.
\end{proof}

\begin{proposition}[machine-verified]\label{prop:continuum}
For every $\lambda\in(1,\varphi)$ there is an explicit injection of
$\{0,1\}^{\mathbb N}$ into the survival itineraries of $1/2$: uncountably many
kind sequences. Moreover $K_\lambda(m)$, the number of kind words of length
$m$, satisfies $K_\lambda(m)\geq\lfloor m/(B+1)\rfloor$ on any window as in
Lemma~\ref{lem:goodchild}.
\end{proposition}

In our normalisation the first assertion recovers, in certified form, a
theorem of Erd\H{o}s, Jo\'o and Komornik \cite{ErdosJooKomornik1990} on the
multiplicity of $q$-expansions for $q<\varphi$; the lemmas above are its
quantitative skeleton, and the count is what the first-moment programme
consumes.

The linear bound is far from the truth: the measured growth of $K_\lambda(m)$
is $(2/\lambda)^m$. It can be replaced by an exponential one, and by an
argument that avoids the return-time analysis altogether.

\begin{proposition}[machine-verified]\label{prop:expcount}
For every $\lambda\in[\tfrac{1000}{667},\tfrac85]$ and every $m$,
\[
  K_\lambda(m)\ \geq\ 2^{\lfloor m/8\rfloor},
\]
and at $\lambda=3/2$, $K_{3/2}(m)\geq15^{\lfloor m/12\rfloor}$, a growth rate
of $15^{1/12}\approx1.2532$ per move against the measured $4/3$.
\end{proposition}

The mechanism is a self-reproducing interval. Say that an interval
$[a,b]$ \emph{doubles in $T$ steps} if from every $x\in[a,b]$ there are two
distinct branch words of length $T$ carrying $x$ back into $[a,b]$; then the
number of surviving words from any point of $[a,b]$ at least doubles every $T$
steps, and iterating gives $2^{\lfloor m/T\rfloor}$. The certificates are
finite: $24$ parameter cells carrying $1333$ interval cells with words of
length at most $8$ for the window, and $503$ cells tiling $[\tfrac14,\tfrac34]$
with $15$ words of length $12$ at $\lambda=3/2$, all checked by kernel
reduction over exact rationals, with the interval arithmetic carried out in the
parameter as well as in the point so that one certificate serves a whole
interval of $\lambda$. Underlying it is a small structural simplification: a
branch word issued from a point of $(0,1)$ is legal exactly when its final
image lies again in $(0,1)$, so a certificate need only control images and
never legality.

Two things this does and does not settle. It settles the counting half of the
first-moment problem, unconditionally and uniformly on a window containing
$3/2$: the supply of kind words grows exponentially, which is the half of the
heuristic of Remark~\ref{rem:candidates} that says one may prepend as many
moves as one likes. It does not settle the other half. Exponentially many kind
words do not have to have well-distributed endpoints, and it is distribution,
not abundance, that Theorem~\ref{thm:density0} requires.

Finally, the two halves meet where they are needed:

\begin{proposition}[common window; machine-verified]\label{prop:commonwindow}
For every $\lambda\in[\tfrac{1000}{667},\tfrac85]$: the transversality of
Proposition~\ref{prop:restricted} applies at $1/\lambda$;
$\lambda^2<\lambda+1$, so Lemmas~\ref{lem:nojump} and~\ref{lem:goodchild}
apply with $\eta=\tfrac15$, $B=3$; the injection of
Proposition~\ref{prop:continuum} exists; and
$K_\lambda(m)\geq\lfloor m/4\rfloor$. The window contains $3/2$.
\end{proposition}

\section{Order, gaps, and a scheduling bound}\label{sec:gaps}

So far we have watched knots one at a time. This section watches the
\emph{spaces between them}, and the payoff is a second proof that Pisot
parameters have bounded knot counts -- shorter than the automaton argument of
Section~\ref{sec:pisot}, and effective: it converts a single number, the
smallest gap in the orbit of $1/2$, directly into a bound on
$N_\lambda$.

Throughout, fix a run and consider the knots alive at a given moment, listed in
increasing order of position. Three elementary lemmas control everything.

\begin{lemma}[order]\label{lem:order}
Every move preserves the spatial order of the knots that survive it.
\end{lemma}

\begin{proof}
The maps $f_0$ and $f_1$ are increasing, so moves $L$ and $R$, which apply one
map to every survivor, preserve order. At a move $M$ the survivors below the
deleted interval are carried by $f_0$ into $(0,1/2)$ and those above it by
$f_1$ into $(1/2,1)$; each group keeps its internal order, and the first group
lands entirely to the left of the second.
\end{proof}

Thus two knots, once created, keep their left-to-right relationship for as long
as both live. Say that a pair of knots \emph{straddles} a move $M$ if one lies
below the deleted interval and the other above it.

\begin{lemma}[gap law]\label{lem:gaplaw}
Let $x<y$ be knots that both survive a move. If the move is $L$, is $R$, or is
an $M$ that the pair does not straddle, then the new distance is
$\lambda(y-x)$. If the pair straddles an $M$, the new distance is
$\lambda(y-x)-(\lambda-1)$.
\end{lemma}

\begin{proof}
In the first case both knots are carried by the same affine map of slope
$\lambda$. In the second, $x\mapsto\lambda x$ and
$y\mapsto\lambda y-(\lambda-1)$.
\end{proof}

So every distance between coexisting knots is stretched by the factor
$\lambda$ at every single move, with one exception: a pair that straddles an
$M$ has its distance reduced by $\lambda-1$ relative to that stretch. Since
all knots lie in $(0,1)$, no distance can ever reach $1$; the stretching must
therefore be paid down, and straddling is the only currency. The final lemma
says the currency is scarce.

\begin{lemma}[one service per move]\label{lem:service}
At any move, among the spatially consecutive pairs of any fixed set of
surviving knots, at most one pair straddles.
\end{lemma}

\begin{proof}
A straddling pair has its two members on opposite sides of the deleted
interval, so a consecutive pair straddles exactly when the deleted interval
falls between its two members -- and an interval falls between the members of
at most one consecutive pair of an ordered list.
\end{proof}

Applied to the extreme pair, the gap law becomes a statement about the whole
configuration. Write $S$ for the \emph{spread}, the distance from the leftmost
to the rightmost knot, and $T=1-S$ for the \emph{deficit}.

\begin{corollary}[deficit law]\label{cor:deficit}
Across a move at which no knot dies and none is born outside the current
range, $T\mapsto\lambda T$ if the extreme pair straddles, and
$T\mapsto\lambda T-(\lambda-1)$ otherwise.
\end{corollary}

\begin{proof}
By Lemma~\ref{lem:gaplaw}, $S\mapsto\lambda S-(\lambda-1)$ in the straddling
case and $S\mapsto\lambda S$ otherwise; substitute $S=1-T$.
\end{proof}

This is verbatim the law of a single knot, with a straddling $M$ in the role of
$R$ and every other move in the role of $L$; and $T$ must remain in $(0,1)$,
since $T\leq0$ says the configuration no longer fits in the strip. A
configuration therefore carries a virtual knot of its own, which has to be kept
alive by the same letters that keep the real ones alive. The tension is visible
in the witnesses: a long run of $L$ and $R$ drives $T$ towards $0$, and only a
straddling $M$ restores it.

These three facts assemble into a pigeonhole argument.

\subsection{Free births}
The three moves are not always three. Relative to a given configuration, an $M$
may act on the knots present exactly as an $L$ or an $R$ would --- same
survivors, same images --- differing only in that it also creates a knot. Call
such an $M$ \emph{free}: it adds a knot at no cost to the configuration.

\begin{lemma}\label{lem:free}
If every knot of $C$ lies below $r/2$ then $M$ acts on $C$ exactly as $R$
does, and if every knot lies above $1-r/2$ then $M$ acts exactly as $L$: in
either case $C$ is contained in one of the two intervals that $M$ leaves
alone, and the $M$ is free. For $\lambda<3/2$ the converses hold as well. For
$\lambda\geq3/2$ they fail: then $r\leq1-r/2$, and a knot sitting in
$[r,1-r/2]$ is deleted by $M$ and by $R$ alike, so $M$ acts on the
one-knot configuration $\{r\}$ exactly as $R$ does although $r\geq r/2$.
\end{lemma}

\begin{proof}
If $C\subset[0,r/2)$ then no knot is deleted by $M$ or by $R$, and both carry
each knot by $f_0$. For the converse suppose $\lambda<3/2$, so $1-r/2<r$, and
let $x\in C$ with $x\geq r/2$. If $x\leq1-r/2$ then $M$ deletes $x$ while $R$,
which deletes only $[r,1]$, does not, since $x\leq1-r/2<r$. If $x>1-r/2$ then
$M$ carries $x$ by $f_1$, whereas $R$ either deletes it or carries it by
$f_0$. Either way $M$ does not act as $R$. The statements for $L$ are the
reflection $x\mapsto1-x$, and the counterexample for $\lambda\geq3/2$ is
checked directly.
\end{proof}

At $\lambda=3/2$ itself the exceptional interval $[r,1-r/2]$ degenerates to the
single point $2/3$, which is not dyadic and so is never occupied --- reachable
positions staying off the cell boundaries is itself certified
(Appendix~\ref{app:verify}); the converse therefore holds for every
configuration that actually arises there.

This gives a clean notion of weak equivalence: two runs are equivalent if,
along the way, each of their moves acts identically on the knots present, so
that an $M$ may be matched with an $L$ or an $R$. Under this equivalence the
positions of the knots depend only on the resulting two-letter word, and the
three-letter structure survives only where some $M$ is not free. Free births
are what one would want in quantity --- a knot gained for nothing --- but they
are also, by Lemma~\ref{lem:free}, births into a configuration huddled in one
of the two end intervals, and such a configuration has small spread and cannot
straddle. So a free $M$ never services a gap, and by Corollary~\ref{cor:deficit} it
moves the deficit as every non-straddling move does --- down toward the fatal
exit at $0$; only the $M$'s that are not free can replenish it. The
economy of the game is visible in the records: in the runs attaining $5$, $6$
and $7$ knots at $\lambda=3/2$, of $48$ moves $M$ only $13$ are free, and the
proportion falls as $k$ grows.

\subsection{The backward picture}
Order preservation has a consequence for the bookkeeping that is worth
isolating, because it removes an apparent exponential.

Read a run backwards. The inverse branches $y\mapsto ry$ and
$y\mapsto ry+(1-r)$ are contractions defined on all of $(0,1)$: backwards
nothing is ever destroyed. A move $L$ applies the second branch to every knot
and $R$ the first; a move $M$ applies the \emph{first} branch to knots below
$1/2$ and the \emph{second} to knots above, since its image is
$[0,r/2)\cup(1-r/2,1]$. So backwards each knot experiences a two-letter
itinerary, and the itineraries of different knots differ only at the moves
$M$. Knots are not destroyed but \emph{annihilated}: a knot reaches exactly
$1/2$ at the move that gave birth to it, and disappears there.

At first sight the moves $M$ impose a compounding constraint --- the letter
that annihilates one knot must serve as $L$ or as $R$ for each older knot,
according to the side that knot happens to occupy, and a knot outliving several
births must fall correctly each time, which looks like $2^{j}$ conditions after
$j$ births. It is not so. By Lemma~\ref{lem:order} the spatial order of the
knots never changes; the knot being annihilated sits at $1/2$; so the knots
below it in the fixed order take one branch and those above it the other. No
choice is made and no condition is imposed. The combinatorial content of a
$k$-knot configuration is therefore a single permutation --- the annihilation
order read against the spatial order --- rather than a tree of side-patterns.

\begin{theorem}[scheduling bound]\label{thm:schedule}
Suppose a run reaches $k$ simultaneous knots, and that throughout the run any
two coexisting knots are at distance at least $\delta>0$. Put
$W=\lceil\log_\lambda(1/\delta)\rceil$. Then
\[
  k\ \leq\ W+\lceil W/2\rceil+1.
\]
\end{theorem}

\begin{proof}
Fix a run with $k$ knots alive at time $n$, and look at the window consisting
of the last $W$ moves. Call a knot \emph{old} if its age is at least $W$; old
knots were alive throughout the window. By Lemma~\ref{lem:order} the old
knots keep a fixed spatial order, so their consecutive pairs are well defined
throughout. Consider one such pair. At time $n-W$ its distance was at least
$\delta$; by Lemma~\ref{lem:gaplaw} that distance is multiplied by $\lambda$
at each of the $W$ moves of the window unless the pair straddles one of them;
and $\lambda^W\delta\geq1$ while distances stay below $1$. So each
consecutive old pair straddles at least one move of the window. By
Lemma~\ref{lem:service} each move is straddled by at most one such pair.
Hence (number of old knots) $-\,1\leq W$.

The remaining knots have age less than $W$, and by the argument of
Proposition~\ref{prop:lowerbound} their birth times are at least two apart, so
there are at most $\lceil W/2\rceil$ of them. Adding the two counts gives the
bound.
\end{proof}

\begin{corollary}[Pisot boundedness, again, and effectively]\label{cor:effective}
Let $\lambda$ be Pisot, let $O$ be the (finite) orbit of $1/2$, and let
$\delta_0$ be the smallest distance between two points of $O$. Then
$N_\lambda\leq W+\lceil W/2\rceil+1$ with
$W=\lceil\log_\lambda(1/\delta_0)\rceil$.
\end{corollary}

\begin{proof}
Every knot is born at $1/2$ and moved by $f_0$ and $f_1$, so all positions lie
in $O$; coexisting knots are distinct by Lemma~\ref{lem:distinct}; so any two
are at distance at least $\delta_0$, and Theorem~\ref{thm:schedule} applies.
\end{proof}

For the golden ratio the orbit is the five-point set of
Theorem~\ref{thm:phi}, whose smallest gap is
$\frac{\varphi-1}{2}-\frac{2-\varphi}{2}=\varphi-\tfrac32=(\sqrt5-2)/2\approx
0.118$; then $W=5$ and $N_\varphi\leq9$. For the plastic number the orbit has
$153$ points and smallest gap at least $239/100000$ (certified; the value is
$\approx0.00239$), giving $W=22$ and $N_\rho\leq34$. Both bounds are now
formally verified (Appendix~\ref{app:verify}). The true values are $2$ and $7$,
so the bound is loose by a factor of about five -- but it needs no reachable
configurations and no automaton, only the orbit and one minimum, and unlike
Theorem~\ref{thm:pisot} it is uniform: the same three lemmas serve every
parameter.

Read in the other direction the theorem says that large knot counts
\emph{force near-collisions}.

\begin{corollary}[near-collisions are necessary]\label{cor:nearcollision}
If some run reaches $k$ simultaneous knots, then for every integer
$m<2(k-2)/3$ there is a moment at which two coexisting knots lie within
$\lambda^{-m}$ of one another.
\end{corollary}

\begin{proof}
If every pair of coexisting knots always kept distance at least
$\lambda^{-m}$, the theorem with $\delta=\lambda^{-m}$, $W=m$ would give
$k\leq m+\lceil m/2\rceil+1\leq \tfrac32m+2$.
\end{proof}

At $\lambda=3/2$ this connects the knot count to the arithmetic of
Section~\ref{sec:threehalves}: positions of knots of age at most $a$ are odd
multiples of $2^{-(a+1)}$, so a pair within $\lambda^{-m}$ is a
near-coincidence among the integers $3^a-\sum\varepsilon_j2^j3^{a-j}$, and how
closely those subset sums can approach one another governs how many knots can
coexist. The erratic behaviour of such approximations is, we believe, the
arithmetic face of the erratic growth of $d_{3/2}(k)$.

Finally, the gap law suggests an accounting heuristic for the quadratic growth
observed at $\lambda=(1+\sqrt2)/2$, which we record without claiming a proof:
with $k$ knots there are $k-1$ consecutive distances, every one of them
stretched by $\lambda$ at every move, and Lemma~\ref{lem:service} allows only
one reduction per move. Each distance must be serviced before it grows to
length $1$, so the servicing load grows linearly with $k$ while the service
rate is fixed; sustaining $k$ knots therefore forces the configuration to
spend an ever-larger share of its moves on maintenance, and the cost of adding
one more knot grows roughly linearly in $k$ -- which integrates to
$d(k)\asymp k^2$. Making this rigorous would require a lower bound on how
small the serviced distances can be made, which is exactly the
anti-concentration question left open above.

\section{Periodic runs and Littlewood polynomials}\label{sec:kind}

A \emph{Littlewood polynomial} is one all of whose coefficients are $+1$ or
$-1$. They enter for a concrete reason: asking whether a knot returns exactly
to its birth point after a fixed number of moves amounts to asking whether a
particular polynomial with coefficients $\pm1$ vanishes at $\lambda$, and
whether such polynomials can vanish at a given number is a much-studied and
difficult question. This section shows the answer here is no, for a reason that
needs none of that theory, and explains why the route the theory would have
supplied is closed. At $\lambda=3/2$ a stronger statement holds --- no knot
ever revisits \emph{any} previous position (Proposition~\ref{prop:norecur}).

Call an infinite sequence $v\in\{L,M,R\}^{\mathbb N}$ \emph{kind} for $\lambda$
if the orbit of $1/2$ under $v$ meets no deleted interval. A knot is immortal
exactly when the sequence of moves following its birth is kind, and if a run has
infinitely many births whose subsequent sequences are kind, and the resulting
knots occupy distinct points, then $N_\lambda$ is unbounded. Periodic runs are
the case that can be treated algebraically, and we fix the indexing as follows:
the run has period $p$, one of its phases carries $M$, and a knot born at that
phase experiences $v=v_1\cdots v_p$ repeated, so that $v_p=M$.

\begin{proposition}\label{prop:littlewood}
The orbit of $1/2$ along a branch itinerary $(\varepsilon_1,\dots,\varepsilon_q)
\in\{0,1\}^q$ satisfies $x_q=1/2$ if and only if
\[
  \sum_{i=0}^{q-1}\bigl(1-2\varepsilon'_i\bigr)\lambda^{\,i}=0,
  \qquad \varepsilon'_i=\varepsilon_{q-i},
\]
that is, if and only if $\lambda$ is a root of the Littlewood polynomial of
degree $q-1$ whose signs record the itinerary.
\end{proposition}

\begin{proof}
The composite map along the itinerary is $x\mapsto\lambda^q
x-(\lambda-1)\sum_{j}\varepsilon_j\lambda^{\,q-j}$. Setting the image of $1/2$
equal to $1/2$ gives $(\lambda^q-1)/2=(\lambda-1)\sum_j\varepsilon_j
\lambda^{\,q-j}$; dividing by $\lambda-1$ and rearranging gives the display,
whose coefficients lie in $\{\pm1\}$.
\end{proof}

For $q=3$ and $\varepsilon=(0,1,1)$ the display reads $1+\lambda+\lambda^2=
2(1+\lambda)$, that is $\lambda^2-\lambda-1=0$, and the orbit of $1/2$ at
$\lambda=\varphi$ is $1/2$, $\varphi/2$, $(3-\varphi)/2$.

Recall from Lemma~\ref{lem:distinct} that under $M$ no surviving knot is carried
to $1/2$: the two preimages of $1/2$ are $r/2$ and $1-r/2$, the endpoints of the
interval $M$ deletes. This alone empties the construction.

\begin{corollary}\label{cor:noperiodic}
Under a periodic run of period $p$ the orbit of a knot born at an $M$ never
returns to $1/2$ at a time divisible by $p$. Equivalently, $1/2$ is not a
periodic point of the associated period map, and
Proposition~\ref{prop:littlewood} has no realisable instances of this kind.
\end{corollary}

\begin{remark}
The corollary rules out one route to unboundedness and not the natural one. It
says that a knot cannot \emph{return to its birth point} under a periodic run;
it says nothing about a knot merely staying alive. Indeed
Proposition~\ref{prop:kindyield} shows that a periodic kind word with an $M$ in
its period would give $N_\lambda=\infty$ immediately, and no such word is
excluded by anything proved here -- which is why the searches reported below,
over all periods up to $14$ at $\lambda=(1+\sqrt2)/2$, were worth running. What
the corollary does exclude is the algebraic shortcut: one cannot exhibit such a
word by solving for a $\lambda$ at which $1/2$ is periodic, because there is
none.
\end{remark}

\begin{proof}
Suppose the orbit returns to $1/2$ at a time $q$ with $p\mid q$. The move at
time $q$ is then the move at phase $p$, namely $M$, and by
Lemma~\ref{lem:distinct} no surviving knot is carried to $1/2$ by $M$.
\end{proof}

\begin{remark}\label{rem:onlymultiples}
The restriction to times divisible by $p$ cannot be removed. The orbit may
return to $1/2$ at other times, since $R$ carries $r/2$ to $1/2$ and $L$ carries
$1-r/2$ to $1/2$, and both of those points survive the move in question. At
$\lambda=\varphi$ the run of period four with $v=RLLM$ carries $1/2$ to
$\varphi/2$, then to $(3-\varphi)/2$, then back to $1/2$ at time three, where
the following $M$ destroys it. Only returns at times divisible by $p$ bear on
periodicity of the period map, and those are what the corollary excludes.
\end{remark}

What survives of the construction is the following, in which the orbit is
allowed to be eventually periodic on a cycle avoiding $1/2$.

Coding a run in base three by $L,M,R\mapsto0,1,2$ realises the kind sequences
as a subset $K_\lambda$ of $[0,1]$, and the threshold of
Section~\ref{sec:thresholds} governs its size.

\begin{proposition}[counts machine-verified]\label{prop:kinddim}
$K_{3/2}$ has exactly $2^n$ cylinders of length $3^{-n}$ at every level $n$;
hence it has Lebesgue measure zero and Hausdorff and box dimension exactly
$\log2/\log3$. It is not self-similar. For $\lambda<3/2$ the dimension is
larger and for $\lambda>3/2$ smaller; at $\varphi$ the set is \emph{sofic} ---
its sequences are exactly the paths through a finite directed graph, here the
five-state one of Theorem~\ref{thm:phi} --- with
$N_{n+3}=4N_n$ cylinders and dimension exactly $2\log2/(3\log3)$.
\end{proposition}

\begin{proof}
At $\lambda=3/2$ the three deleted intervals tile, so a point lies in exactly
one of them: exactly one letter is fatal at each node of the survival tree of
$1/2$, never none and never two, and the cylinder count doubles at each level.
Measure $(2/3)^n\to0$; and a set covered at level $n$ by $2^n$ intervals of
length $3^{-n}$, each containing a further such set, has both dimensions
$\log2/\log3$. Self-similarity fails because which of the three subintervals is
deleted depends on the current position, which varies from node to node. For
$\lambda<3/2$ the deleted intervals leave gaps, so some nodes have three
children; for $\lambda>3/2$ they overlap, so some have one. At $\varphi$ the
survival tree is generated by the five-state automaton of
Theorem~\ref{thm:phi}, in which two of the states have two children and three
have one; the counts satisfy $N_{n+3}=4N_n$, giving growth $4^{1/3}$ per level.
\end{proof}

\begin{figure}[ht]
\centering
\includegraphics{kindsets.pdf}
\caption{The kind sequences, coded in base three by $L,M,R\mapsto0,1,2$, drawn
to level six with the cylinder counts at the right. The middle third is absent
from every picture: a run beginning with $M$ kills the newborn at once. At
$3/2$ exactly one letter is fatal at each node, so exactly one third is deleted
at each step --- the Cantor count $2^n$, though the third that goes varies from
node to node, so the set is not self-similar. At $\sqrt2$ the deleted intervals
leave gaps, so some nodes keep all three children and the count breaks away
from $2^n$ at level six. At $\varphi$ they overlap, some nodes keep only one,
and the count stalls every third level.}
\label{fig:kindsets}
\end{figure}

The measured dimensions are $0.834$ at $(1+\sqrt2)/2$, $0.679$ at $\sqrt2$ and
$0.421\approx2\log2/(3\log3)$ at $\varphi$, bracketing the exact
$\log2/\log3=0.6309$ at $3/2$ as the proposition requires.

\begin{proposition}[machine-verified]\label{prop:kindyield}
If $v$ is kind, with $v_p=M$, then $N_\lambda=\infty$.
\end{proposition}

\begin{proof}
The knots born at the successive occurrences of that phase are all immortal, and
by Lemma~\ref{lem:distinct} each is distinct from all its predecessors at the
moment of its birth and remains so. Hence the number of knots present grows
without bound.
\end{proof}

No hypothesis on the orbit is needed. Coincidence of two immortal knots would
require the orbit of $1/2$ to repeat a value, and Lemma~\ref{lem:distinct}
forbids two knots from ever meeting.

\begin{remark}
Neither $\sqrt2$ nor $3/2$ is a root of a Littlewood polynomial, so at those two
parameters $1/2$ is not a periodic point for any itinerary whatever, not merely
for those ending in $M$. For $\sqrt2$: separating a Littlewood polynomial
evaluated at $\sqrt2$ into rational and irrational parts gives $A+B\sqrt2$ with
$A$ a $\{\pm1\}$-combination of distinct powers of $2$, whose leading term
exceeds the sum of the others, so $A\neq0$. For $3/2$: clearing denominators
gives $\sum_k a_k2^k3^{\,d-k}=0$, in which only the term $k=d$ is odd.
\end{remark}

A search over all itineraries with $q\leq10$, conducted in exact arithmetic in
$\Q(\lambda)$, confirms Corollary~\ref{cor:noperiodic}. The numbers of
itineraries whose orbit is realisable, that is for which every step satisfies
the appropriate strict inequality, are $0,2,2,10,10,44,56,180,250$ for
$q=2,\dots,10$; none of them admits an assignment of moves periodic in any
$p\mid q$ with $M$ at phase $p$.

Two further computations bear on the remaining possibility, that of a kind
periodic run whose orbit is not eventually periodic. For the sequence
$v=RM$ the set of $\lambda\in(1,2)$ for which $1/2$ survives $n$ moves is
contained in $(1,1.2)$ for $n\geq25$, and its measure decays roughly as
$n^{-0.83}$, so short periods serve only parameters near $1$. At
$\lambda=\sqrt2$ and $\lambda=3/2$, over all periodic sequences of period at
most $9$ containing an $M$, the point $1/2$ survives at most $22$ and $17$ moves
respectively.



\section{A density criterion for unboundedness}\label{sec:density}

\emph{Relaxation in force: a hypothesis is assumed, and the conclusion sought
is strengthened. The hypothesis is proved for no non-Pisot parameter.}

This section isolates a single hypothesis about the orbit of $1/2$ -- a
quantitative density statement -- and shows that it implies
$N_\lambda=\infty$. The hypothesis is measured to hold, at the best possible
rate, at every non-Pisot parameter we can test; and it fails, provably, at
every Pisot parameter. We believe it is the mechanism of the dichotomy.

Two preparatory lemmas. Recall from Theorem~\ref{thm:mono} that prepending a
letter never destroys a knot; the first lemma sharpens this: prepending does
not even move one.

\begin{lemma}\label{lem:permanent}
Let $v$ be a word and $c$ a letter. The knots of $cv$ are exactly the knots of
$v$, with the same positions and ages, together with one new knot of age
$|v|$, present exactly when $c=M$ and $1/2$ survives $v$.
\end{lemma}

\begin{proof}
A knot alive at the end of a word is determined by the suffix following its
birth (Lemma~\ref{lem:suffix}): it exists iff its birth letter is $M$ and
$1/2$ survives that suffix, and its position is the image of $1/2$ under that
suffix. Prepending $c$ leaves every proper suffix of $cv$ that lies inside $v$
unchanged, and adds one new birth position at the front.
\end{proof}

So a run extended into the past accumulates knots and never disturbs the ones
it has; the only question is whether it can keep gaining. Write $S(v)$ for the
set of starting points in $(0,1)$ that survive the word $v$.

\begin{lemma}\label{lem:survivorset}
$S(v)$ is a union of at most $1+\#\{M\text{'s in }v\}$ open intervals, of
total length exactly $r^{|v|}$; in particular it contains an open interval of
length at least $r^{|v|}/(|v|+1)$.
\end{lemma}

\begin{proof}
Induct from the right: $S(\varnothing)=(0,1)$, and $S(cv')$ is the image of
$S(v')$ under the inverse branch of the move $c$, which is affine of slope $r$
for $c=L,R$ and piecewise affine of slope $r$ with two order-preserving
branches for $c=M$. Each step multiplies total length by $r$ and increases the
number of intervals by at most one, at an $M$; and the intervals are open,
survival being defined by strict inequalities.
\end{proof}

By Lemma~\ref{lem:permanent}, the extended word $Muv$ has one more knot than
$v$ exactly when $1/2$ survives $uv$ -- that is, when the image of $1/2$ under
$u$ lands in $S(v)$ while surviving $u$. Gaining is thus a question about how
densely the kind orbit of $1/2$ fills the interval. This motivates:

\begin{quote}
\textbf{Hypothesis (D$_\lambda$).} There is a constant $C$ such that every
interval $J\subseteq(0,1)$ of length $\varepsilon$ contains the endpoint of a
kind word of length at most $C\log(1/\varepsilon)$.
\end{quote}

\begin{theorem}\label{thm:density}
If \textup{(D$_\lambda$)} holds then $N_\lambda=\infty$. Moreover the witnesses
form a single nested family: each is obtained from the last by prepending, so
every knot, once created, appears with the same age in all later witnesses,
and $d_\lambda(k)$ grows at most exponentially in $k$.
\end{theorem}

\begin{proof}
Given a word $v_k$ carrying $k$ knots, Lemma~\ref{lem:survivorset} provides a
component of $S(v_k)$ of length at least $r^{|v_k|}/(|v_k|+1)$, and
(D$_\lambda$) a kind word $u$ of length at most
$C\bigl(|v_k|\log\lambda+\log(|v_k|+1)\bigr)$ whose endpoint lies in it. Then
$1/2$ survives $uv_k$, and $v_{k+1}=Muv_k$ carries $k+1$ knots by
Lemma~\ref{lem:permanent}. Iterating from $v_1=M$ gives
$|v_{k+1}|\leq(1+C\log\lambda)|v_k|+O(\log|v_k|)$, and
$N_\lambda(|v_k|)\geq k$ for every $k$.
\end{proof}

Before stating the criterion we record the form the question takes when the
run is read backwards, which is where the construction above actually lives.
Reading a run from its last move, each knot moves by one of the two inverse
branches $y\mapsto ry$ and $y\mapsto ry+(1-r)$ (Section~\ref{sec:gaps}); these
are defined on all of $(0,1)$, so backwards nothing is ever destroyed. What
happens instead is that a knot arrives at exactly $1/2$ at the move that gave
birth to it, and is \emph{annihilated} there. The backward game is thus a
pursuit of a single target by a family of points driven by one control.

\begin{lemma}\label{lem:oneperstep}
At each backward step at most one point is annihilated.
\end{lemma}

\begin{proof}
The two inverse branches have disjoint images, $[0,r)$ and $[1-r,1]$ at a move
$L$ or $R$ and $[0,r/2)$ and $(1-r/2,1]$ at a move $M$, and each is injective;
so every composite of inverse branches is injective, and the point carried to
$1/2$ after $a$ steps, if there is one, is unique.
\end{proof}

Consequently the set of points ever annihilated during a backward play is at
most countable, whatever the run --- so \emph{countably infinite} is the
largest it can be, and asking for infinitely many annihilations is asking for
the most that can be asked. That question is the dual of the one this paper
began with, and it is not the same question: unboundedness says the finite
counts $N_\lambda(n)$ grow without bound, while this asks for one play in which
infinitely many points are eventually annihilated. The second implies the
first; whether the first implies the second is exactly the gap discussed below.
Note also that a positive answer does not make any configuration infinite in
the forward game, nor does it exhaust the annihilated set in the backward one:
by Lemma~\ref{lem:oneperstep} only finitely many points are gone at any finite
stage.

Lemma~\ref{lem:permanent} has a second consequence, which converts the
existence of an \emph{infinite} configuration into a statement about finite
witnesses. Given a word $w$ carrying $k$ knots, let its \emph{age vector} be
the increasing list $a_1<\dots<a_k$ of the ages of those knots.

\begin{theorem}[compactness criterion]\label{thm:compactness}
The following are equivalent.
\begin{enumerate}
\item Some left-infinite run carries infinitely many simultaneous knots.
\item In the backward reading of some run, infinitely many points are
  annihilated --- equivalently, by Lemma~\ref{lem:oneperstep}, the set of
  annihilated points is countably infinite.
\item There are words $w_k$ carrying $k$ knots whose age vectors are
  pointwise bounded: for each $i$, $\sup_k a_i^{(k)}<\infty$.
\end{enumerate}
Any of them implies $N_\lambda=\infty$.
\end{theorem}

\begin{proof}
Statements (i) and (ii) are the same statement read in the two directions: a
knot of age $a$ alive at the end is a point annihilated at backward step $a$,
and conversely.

If $\omega$ is such a run, its truncations after $n$ moves carry arbitrarily
many knots, and the $i$th smallest age is the same in all of them once it
appears, by Lemma~\ref{lem:permanent}; so (i) implies (iii).

Conversely, read each $w_k$ backwards, as an element of
$\{L,M,R\}^{\mathbb N}$ padded arbitrarily. That space is compact in the
product topology, so some subsequence converges; convergence means agreement on
longer and longer \emph{suffixes} of the original words. By
Lemma~\ref{lem:suffix} the presence of a knot of age $a$ depends only on the
last $a+1$ letters. Fix $i$ and let $A_i=\sup_ka_i^{(k)}<\infty$. Along the
subsequence the values $a_i^{(k)}\leq A_i$ are integers, so after passing to a
further subsequence they are constant, $=\alpha_i$; and the last $\alpha_i+1$
letters are eventually constant too. Hence the limit word has a knot of age
$\alpha_i$ for every $i$, and the $\alpha_i$ are distinct. A diagonal
subsequence serves all $i$ at once.
\end{proof}

The criterion is exact for the stronger statement, an infinite configuration,
which implies $N_\lambda=\infty$; whether $N_\lambda=\infty$ implies it back is
not known. Its hypothesis is checkable: at $\lambda=(1+\sqrt2)/2$ a search
minimising the age vector rather than the length yields
\[
\begin{array}{r@{\;}l}
k=7:&\ 0,2,4,6,8,10,12\\
k=8:&\ 0,2,4,6,8,10,12,18\\
&\ \vdots\\
k=13:&\ 0,2,4,6,8,10,12,18,24,26,31,39,48\\
k=14:&\ 0,2,4,6,8,10,12,18,24,26,31,39,48,62
\end{array}
\]
each vector extending the previous one, so that $a_i$ is not merely bounded but
constant in $k$ as far as computed. By Lemma~\ref{lem:permanent} the nesting is
automatic once the witnesses are built by prepending; what is not automatic,
and not known, is that the process continues. The $k=13$ entry is the source of
the bound $d_{(1+\sqrt2)/2}(13)\leq49$ in Section~\ref{sec:nonpisot}.

In fact, for the bare conclusion $N_\lambda=\infty$ the rate in
(D$_\lambda$) can be dispensed with entirely.

\begin{theorem}[density criterion, topological form]\label{thm:density0}
If the endpoints of kind words are dense in $(0,1)$, then some run carries
infinitely many simultaneous knots; in particular $N_\lambda=\infty$.
\end{theorem}

\begin{proof}
Given a word $v$ carrying $k$ knots, Lemma~\ref{lem:survivorset} provides an
open interval $J\subseteq S(v)$. Density provides a kind word $u$ with
$\Phi_u(1/2)\in J$; then $1/2$ survives $uv$, and $Muv$ carries $k+1$ knots
with every previous knot's position and age unchanged
(Lemma~\ref{lem:permanent}). Iterating from $v_1=M$ gives words $v_k$ with
$k$ knots whose age vectors are pointwise eventually constant, and
Theorem~\ref{thm:compactness} supplies the single run.
\end{proof}

\begin{remark}[the candidate count]\label{rem:candidates}
The theorem makes precise a counting heuristic for the backward game which is
worth recording, because it locates exactly what density is doing. Suppose a
backward play over the word $w$ already annihilates $m$ knots, and one wants
one more. A \emph{candidate} is an intended itinerary through $w$: at each $M$
of $w$, a declaration of which branch the new point will take. The entry
points that realise a given itinerary and survive form an open interval ---
empty for most declarations, which are bogus in the sense that any point
following them would land on the wrong side of some $M$ --- and the nonempty
cells partition $S(w)$, so their lengths sum to exactly $r^{|w|}$. For the
$19$-move run attaining five knots at $\lambda=3/2$ there are $2^{11}=2048$
candidates, of which six are genuine, tiling a set of length exactly
$(2/3)^{19}$. The decisive asymmetry is that this target is \emph{fixed},
while the supply of entry points reachable from $1/2$ grows without bound with
the number $p$ of further moves one is prepared to prepend --- at the measured
rate $(2/\lambda)^p$ --- and $p$ is free. Density is the statement that an
unbounded supply eventually meets any fixed open target; Hypothesis
(D$_\lambda$) additionally bounds how soon, which is what the exponential
bound on $d_\lambda(k)$ in Theorem~\ref{thm:density} costs.
\end{remark}

Both criteria are certified: the topological form as stated, and the
quantitative form with the explicit bound $d_\lambda(k)\leq B^k-1$ for
$B=\max\bigl(2,\,1+C(\log\lambda+1)\bigr)$, each carrying its unproved
hypothesis visibly in its statement (Appendix~\ref{app:verify}).

Both hypotheses fail at every Pisot $\lambda$, since the kind orbit
of $1/2$ is then a finite set and misses every sufficiently short interval in
its gaps -- as Theorem~\ref{thm:pisot} demands it must, given the theorem
above. Whether it holds at any specific non-Pisot parameter is open, and at
$\lambda=3/2$ it sits in distinguished company: the endpoint of a kind word of
length $a$ is an odd multiple of $2^{-(a+1)}$ determined by the subset sums of
Section~\ref{sec:threehalves}, and density statements for such orbits are
adjacent to the distribution of $(3/2)^n$ modulo one and to Mahler's
$Z$-numbers \cite{Mahler1968}. For the \emph{unconstrained} words the
endpoints equidistribute exactly: the map
$(\varepsilon_j)_{j<\ell}\mapsto\sum_{j<\ell}\varepsilon_j2^j3^{a-j}\bmod
2^\ell$ is a bijection, each term being $2^j$ times an odd number. The entire
difficulty is the survival constraint.

What can be measured supports the hypothesis at the optimal rate. At
$\lambda=3/2$, in exact arithmetic, let $m^*(L)$ be the least depth by which
the kind orbit has visited every dyadic interval of length $2^{-L}$; counting
alone forbids this before $L\log2/\log(4/3)\approx2.41\,L$, and the measured
values $m^*(4),\dots,m^*(9)=6,11,13,16,19,22$ approach that bound from below
with ratios $1.50,2.20,2.17,2.29,2.38,2.44$: the orbit spreads essentially as
fast as its cardinality permits. A floating-point survey across eighteen
parameters in $(1,2)$ shows the same counting-rate coverage at every non-Pisot
point tested, with failure occurring exactly at the four Pisot points on the
grid, by orbit closure (Appendix~\ref{app:verify}).

Finally, the metric precedents. For almost every base $\beta>1$ the sequence
$(\beta^n)$ equidistributes modulo one (Koksma \cite{Koksma1935}); for almost
every $\lambda\in(1,2)$ the Bernoulli convolution with ratio $1/\lambda$ --
which is precisely the law of the backward orbit here -- is absolutely
continuous (Solomyak \cite{Solomyak1995}). Both are statements that most
intervals receive their share; (D$_\lambda$) asks for all intervals, and the
construction of Theorem~\ref{thm:density} in fact needs less: not every
interval, only the particular ones it asks about, chosen one at a time as the
construction proceeds. For an almost-everywhere theorem we expect the right tool to be
\emph{parameter exclusion}: build one construction that works for a whole
interval of $\lambda$ at once, and at each stage discard the small set of
$\lambda$ for which that stage fails, arranging the discarded sets to have
summable measure so that most $\lambda$ survive every stage. This is the
strategy Benedicks and Carleson introduced for a different family, and the
estimate it needs here --- that each stage fails only on a short range of
$\lambda$ --- is what Section~\ref{sec:transversality} supplies. We pose the
question below.

\section{The conjugate modulus}\label{sec:conjugate}

Among Pisot parameters the game is bounded, but the bounds differ widely: $2$
at the golden ratio and $7$ at the plastic number. This section identifies the
quantity responsible. It is not $\lambda$, and it is not the threshold of
Section~\ref{sec:thresholds}; it is the modulus of the conjugates.

Recall from Remark~\ref{rem:overlap} that the backward form of the game is the
iterated function system $\{rx,\ rx+(1-r)\}$, whose invariant measure is the
Bernoulli convolution $\nu_r$. As a convolution of two-point masses its Fourier
transform is an explicit product of cosines,
\[
  |\widehat{\nu_r}(\xi)|=\prod_{j\geq0}\bigl|\cos\bigl(\pi(1-r)r^j\xi\bigr)\bigr| ,
\]
which at $r=\tfrac12$ telescopes to $|\sin\pi\xi/\pi\xi|$, the transform of
Lebesgue measure -- as it must, since there the two branches tile. Evaluating
along the frequencies $\xi_N=\lambda^N/(1-r)$ turns every argument into
$\pi\lambda^{N-j}$, so
\[
  |\widehat{\nu_r}(\xi_N)|=\prod_{m\leq N}\bigl|\cos(\pi\lambda^m)\bigr|
  \ \xrightarrow[N\to\infty]{}\ \prod_{m\in\Z}\bigl|\cos(\pi\lambda^m)\bigr| .
\]
For Pisot $\lambda$ this limit is positive: the factors with $m\to-\infty$
tend to $1$ because $\lambda^m\to0$, and those with $m\to+\infty$ tend to $\pm1$
because $\|\lambda^m\|\to0$. That positivity is Erd\H{o}s's theorem
\cite{Erdos1939}; what the product adds is a \emph{number}, and the number
turns out to be informative. Write $\|\lambda^m\|$ for the distance to the
nearest integer and $c(\lambda)$ for the maximum modulus of the conjugates of
$\lambda$, so that $\|\lambda^m\|$ decays like $c(\lambda)^m$.

\begin{center}
\begin{tabular}{lccccc}
\toprule
$\lambda$ & $c(\lambda)$ & $|\Orb(\lambda)|$ & $\delta_0$ &
$\prod_{m\in\Z}|\cos(\pi\lambda^m)|$ & $\max N_\lambda$\\
\midrule
golden $1.61803$ & $0.61803$ & $5$ & $0.118$ & $6.6135\times10^{-3}$ & $2$\\
tribonacci $1.83929$ & $0.73735$ & $7$ & $0.0674$ & $1.2475\times10^{-3}$ & $3$\\
supergolden $1.46557$ & $0.82603$ & $43$ & $0.0109$ & $6.6208\times10^{-6}$ & $4$\\
plastic $1.32472$ & $0.86884$ & $153$ & $0.00239$ & $3.3942\times10^{-7}$ & $7$\\
\bottomrule
\end{tabular}
\end{center}

Every column is monotone in $c(\lambda)$, and none is monotone in $\lambda$:
the tribonacci constant is the largest of the four parameters and has the
second smallest maximum. The mechanism is the same on both sides of the table.
A conjugate that contracts slowly means that $\|\lambda^m\|$ decays slowly, so
that many cosine factors sit well away from $\pm1$ before the product
stabilises -- hence a low Fourier floor; and it means that the orbit of $1/2$
takes longer to close, so it has more points packed into $(0,1)$ -- hence a
small minimum gap $\delta_0$. The Fourier floor and the minimum gap are two
measurements of one quantity, in frequency and in space.

Only the third link is a theorem. Theorem~\ref{thm:schedule} converts
$\delta_0$ into a genuine bound on $N_\lambda$, and it is what produces the
values $9$ and $34$ of Corollary~\ref{cor:effective}. The correlation of the
last column with the rest of the table is an observation on four parameters,
not a proof, and we do not know a bound on $N_\lambda$ in terms of
$c(\lambda)$ directly. What the table does establish is that the size of the
bound is a question about conjugates rather than about $\lambda$ itself.

\begin{figure}[ht]
\centering
\includegraphics{densities.pdf}
\caption{The limit distribution of the backward set --- the Bernoulli
convolution $\nu_r$ --- computed by iterating the transfer operator, at
resolution $300$. Below each panel is the mean absolute difference between
adjacent bins, relative to the mean, at resolutions $3000$ and $12000$. Only
the behaviour under refinement is diagnostic: at $\varphi$ it does not move,
the signature of a singular measure; at the plastic number it falls and then
stops, so that a profile which looks smooth to the eye is singular as well; at
half the silver ratio it continues to fall, as an actual density must.}
\label{fig:densities}
\end{figure}

\begin{remark}[the plastic number is isolated]\label{rem:isolated}
By a theorem of Siegel the plastic number is the smallest Pisot number, so the
interval $(1,\rho)$ contains none at all; by a theorem of Salem the smallest
limit point of the Pisot set is $\varphi$, so below $\varphi$ the Pisot numbers
are isolated, the next one after $\rho$ being the root of $x^4=x^3+1$ at
$1.3803$. Two consequences. First, if boundedness occurs exactly at Pisot
parameters, then $N_\lambda=\infty$ for \emph{every} $\lambda$ in $(1,\rho)$,
an interval containing $(1+\sqrt2)/2$ and $5/4$. Second, the singularity at
$\rho$ is a solitary phenomenon, and one can watch it dissolve: computing the
local variation of the transfer-operator profile at $\rho+\varepsilon$ over a
hundredfold refinement, the value at $\varepsilon=0$ settles at $0.0027$ and
does not move, while at $\varepsilon=3\times10^{-4}$ it continues to fall, and
by $\varepsilon=10^{-2}$ it falls by the factor $4$ per fourfold refinement
characteristic of an absolutely continuous density. At the coarsest resolution
all of these agree to two digits.
\end{remark}

\begin{remark}
These products demand care. Evaluating
$\prod_{|m|\leq M}|\cos(\pi\lambda^m)|$ in $60$-digit arithmetic gives
$2.6\times10^{-10}$ at the tribonacci constant, wrong by five orders of
magnitude, because $\lambda^{200}$ then exceeds the working precision and its
fractional part is noise; the same computation at $400$ digits is stable in $M$
and agrees with the transform to five figures. This is the phenomenon of
Appendix~\ref{app:verify} in another guise: the quantities that decide the
answer here are exactly the ones floating-point arithmetic destroys.
\end{remark}

\begin{figure}[ht]
\centering
\includegraphics{records32.pdf}
\caption{Record configurations at $\lambda=3/2$: for each $k$, the knots
present after $d_{3/2}(k)$ moves, at the first moment $k$ knots coexist. The
dashed line marks the birth point $1/2$, which is occupied in every record,
since the last move of a shortest run must be the birth of the $k$th knot.
The particular configurations drawn nest along the two chains
$k=2\subset3\subset5\subset7$ and $k=4\subset6$; the general fact behind this
is Proposition~\ref{prop:twostep}. From $k=4$ onward the smallest distance
between two knots is $13/512$ at every record, realised by the newborn at
$1/2$ together with a knot at $243/512$ or at $269/512$; it does not shrink as
$k$ grows to $7$, consistent with Corollary~\ref{cor:nearcollision}, which at
$k=7$ demands only a gap below $(3/2)^{-3}$.}
\label{fig:records32}
\end{figure}

\begin{proposition}[observed]\label{prop:twostep}
At $\lambda=3/2$, for $1\leq k\leq5$, every configuration attaining $k+2$
knots in $d_{3/2}(k+2)$ moves contains, as a set of points, some
configuration attaining $k$ knots in $d_{3/2}(k)$ moves. The corresponding
statement one step apart is false: of the four configurations attaining five
knots at depth $19$, only two contain a four-knot record.
\end{proposition}

This is a computation, not a theorem: the record configurations were
enumerated exhaustively at each depth --- there are $2,2,2,4,2$ of them for
$k=2,\dots,6$, and $4$ at $k=7$ --- and every one of the $2,2,4,2,4$
configurations at levels $3,4,5,6,7$ contains a record two levels below. The
obvious mechanism does not explain it. Prepending preserves knots and their
positions exactly (Lemma~\ref{lem:permanent}), so containment would follow if
the shorter record word were a suffix of the longer; it is not, already at
$k=5$ and $k=7$, where the point sets nest but the words are unrelated. We do
not know why two steps is the right distance, nor whether the pattern persists.

\begin{question}
Is Proposition~\ref{prop:twostep} a theorem, and is the step of two essential?
\end{question}

\section{\texorpdfstring{The case $\lambda=3/2$}{The case lambda = 3/2}}\label{sec:threehalves}

\emph{Relaxation in force: reduction to a known open problem. What follows
explains the difficulty and does not remove it.}

At $\lambda=3/2$ the three deleted intervals are the three ternary cells of
$[0,1]$, so the rules can be stated in base three. Knot positions are dyadic, by
the computation in Section~\ref{sec:nonpisot}, hence never equal to $1/3$ or
$2/3$, and the leading ternary digit
\[
  D(x)=\lfloor 3x\rfloor\in\{0,1,2\}
\]
is defined at every knot.

\begin{proposition}[machine-verified]\label{prop:ternary}
Identify the moves $L,M,R$ with the digits $c=0,1,2$. Then a move consists of
choosing a forbidden digit $c$; every knot with $D(x)=c$ is destroyed; every
surviving knot moves to $\bigl(3x-[\,D(x)>c\,]\bigr)/2$; and when $c=1$ the
point $1/2$ is adjoined.
\end{proposition}

\begin{proof}
The interval deleted by $L$ is $[0,1/3]$, by $M$ is $[1/3,2/3]$ and by $R$ is
$[2/3,1]$, which are the cells of $D=0,1,2$. A survivor lies below the deleted
cell exactly when $D(x)<c$, and is then moved by $f_0(x)=3x/2$; otherwise it
lies above and is moved by $f_1(x)=(3x-1)/2$.
\end{proof}

Three consequences are immediate at this parameter. Since $D(1/2)=1$, a newly
created knot lies in the very cell that the creating move forbids, which is
Lemma~\ref{lem:distinct} here. A configuration admits no move preserving all of
its knots exactly when it meets all three cells, so two knots can always be
preserved and three is the first count at which obstruction arises. And writing
$e=D-[\,D>c\,]$ for the surviving digit's rank among the two survivors,

\begin{center}
\begin{tabular}{lccc}
\toprule
 & $c=0$ & $c=1$ & $c=2$\\
\midrule
$D=0$ & --- & $0$ & $0$\\
$D=1$ & $0$ & --- & $1$\\
$D=2$ & $1$ & $1$ & ---\\
\bottomrule
\end{tabular}
\end{center}

so a knot in an outer cell moves by a branch determined by its own digit alone.
Only a knot in the middle cell is sensitive to which of the two sparing moves is
played, and newly created knots are exactly the ones in the middle cell.

\subsection{Three coordinates}

\begin{proposition}[machine-verified]\label{prop:base32}
At $\lambda=3/2$ Proposition~\ref{prop:itinerary} reads
$\sum_{j\geq1}\varepsilon_j(2/3)^{j}=1$: the itinerary of a knot is a
$\{0,1\}$-representation of $1$ in base $3/2$. Moreover
$D(x_n)=\varepsilon_{n+1}+[\,2x_{n+1}>1\,]$.
\end{proposition}

\begin{proof}
The first claim is Proposition~\ref{prop:itinerary} with
$1/(2(\lambda-1))=1$. For the second, $3x_n=\varepsilon_{n+1}+2x_{n+1}$ and
$2x_{n+1}<2$.
\end{proof}

Clearing denominators in the truncation gives the integer form. Writing
\[
  m_n=3^n-\sum_{j=1}^{n}\varepsilon_j2^{j}3^{\,n-j},
\]
one has $x_n=m_n/2^{\,n+1}$ with $m_n$ odd and $0<m_n<2^{\,n+1}$, and
$D(x_n)$ equals $0$, $1$ or $2$ according as $3m_n$ lies below $2^{\,n+1}$,
between $2^{\,n+1}$ and $2^{\,n+2}$, or above $2^{\,n+2}$. So the rule is
integer arithmetic on powers of three read in base two: triple, compare against
the two current powers of two, and subtract one of them.

\begin{proposition}[machine-verified]\label{prop:mahler}
Set $w=2x\in(0,2)$ and write $w=p+y$ with $p=\lfloor w\rfloor\in\{0,1\}$ and
$y\in[0,1)$. Then
\[
  D=\lfloor \tfrac32 w\rfloor,\qquad
  w'=\tfrac32w-\varepsilon,\qquad
  y'=\bigl\{\tfrac32y+\tfrac p2\bigr\},\qquad
  D=p+\bigl[\,y\geq\tfrac{2-p}{3}\,\bigr],\qquad p'=D-\varepsilon .
\]
\end{proposition}

The recursion for $y$ is that of Mahler \cite{Mahler1968}: if
$\xi(3/2)^n=A_n+y_n$ then $y_{n+1}=\{\tfrac32y_n+\tfrac12[A_n\text{ odd}]\}$,
and the coordinate $p$ plays the role of the parity of $A_n$. The systems differ
in two respects. Mahler constrains a single orbit to $y<1/2$ for all $n$,
whereas here the window is all of $(0,2)$ and no single knot is constrained at
all; and here the carry $p$ is not determined by a real number but chosen,
subject to the digit condition, by a control sequence shared among all the
knots. The difficulty has moved from the width of the window to the sharing of
the control.

We record a consequence of the dyadic structure which strengthens the
no-return theorem of Section~\ref{sec:kind} from one point to all of them.

\begin{proposition}\label{prop:norecur}
At $\lambda=3/2$ no knot ever occupies the same position twice.
\end{proposition}

\begin{proof}
Every position is a dyadic rational, the game acting on numerators over the
denominator $2^{n+1}$. If a knot's position $x$ recurred after $B$ further
moves, then $x$ would be a fixed point of the affine map those $B$ moves apply
to it: $x\mapsto\lambda^Bx-c$ with
$c=(\lambda-1)\sum_{j\in S}\lambda^{B-j}$ for the set $S$ of steps at which
the knot took the upper branch. Solving, and clearing powers of $2$,
\[
  x\;=\;\frac{\sum_{j\in S}3^{B-j}2^{j-1}}{3^B-2^B}.
\]
The denominator is odd, so $x$ is dyadic only if it divides the numerator. But
the numerator lies between $0$ and $\sum_{j=1}^B3^{B-j}2^{j-1}=3^B-2^B$, with
the extremes attained only at $S=\varnothing$ and $S=\{1,\dots,B\}$; so the
only dyadic fixed points are $0$ and $1$, and no knot is ever there.
\end{proof}

This rules out a second candidate mechanism for
Proposition~\ref{prop:twostep}: the two-step containments cannot arise from a
block of moves acting as the identity on the old knots, since an identity
block would make each of their positions a fixed point of its composite. The
positions in the larger record are the same \emph{numbers}, but every knot
carrying them has arrived by a path that never pauses.

\subsection{A knot-free equivalent}

\begin{proposition}[machine-verified]\label{prop:translation}
$N_{3/2}$ is unbounded if and only if for every $k$ there are an $N$, a sequence
$c\in\{0,1,2\}^N$ and indices $t_1<\dots<t_k$ with $c_{t_i}=1$, such that each
of the $k$ trajectories
\[
  w_{t_i}=1,\qquad
  w_{s+1}=\tfrac32w_s-\bigl[\,\lfloor\tfrac32w_s\rfloor>c_{s+1}\,\bigr]
\]
satisfies $\lfloor\tfrac32w_s\rfloor\neq c_{s+1}$ for all $t_i\leq s<N$.
\end{proposition}

This is the strongest sense in which the reduction is exact: not an analogy
and not one direction of an implication, but an equivalence, and the
equivalence itself has been checked by machine
(\texttt{unbounded\_iff\_mahler}). Whoever settles the displayed question about
the recursion, in either direction, settles Lawler's problem.

\begin{proof}
Proposition~\ref{prop:ternary} and Proposition~\ref{prop:mahler} identify the
game with this recursion, and Lemma~\ref{lem:suffix} identifies the knot count
with the number of births whose trajectories run to completion.
\end{proof}

\begin{proposition}\label{prop:immortal32}
If some $c\in\{0,1,2\}^{\mathbb N}$ has infinitely many indices $t$ with $c_t=1$
whose trajectories never meet the control, then $N_{3/2}=\infty$.
\end{proposition}

\begin{proof}
Those knots are immortal, and distinct by Lemma~\ref{lem:distinct}.
\end{proof}

\begin{remark}\label{rem:converse}
Whether the converse holds is open. Unboundedness supplies a finite witness for
each $k$, and prepending never destroys a knot, so the witnesses assemble into a
tree, finitely branching, on which the knot count is monotone away from the root
and unbounded overall. That does not force a single branch along which the count
is unbounded, and Section~\ref{sec:sqrt2} exhibits configurations from which no
extension gains anything. Unboundedness through knots of ever longer but finite
lifetime is not excluded.
\end{remark}

\section{\texorpdfstring{The case $\lambda=\sqrt2$}{The case lambda = sqrt 2}}\label{sec:sqrt2}

For $\lambda=\sqrt2$ the positions lie in $\tfrac12\Z[\sqrt2]$ and the
denominators do not grow: writing $x=(a+b\sqrt2)/2$ with $a,b\in\Z$, one has
$f_0(a,b)=(2b,a)$ and $f_1(a,b)=(2b+2,a-2)$, and a new knot is $(1,0)$.

\begin{proposition}\label{prop:blocks}
Composing $f_\varepsilon$ followed by $f_\delta$ gives
\[
  x\ \longmapsto\ 2x+(\delta-2\varepsilon)+(\varepsilon-\delta)\sqrt2 .
\]
In particular the block $RR$ preserves exactly the knots with $x<1/2$ and acts
on them as $x\mapsto2x$, while the block $LL$ preserves exactly the knots with
$x>1/2$ and acts on them as $x\mapsto2x-1$.
\end{proposition}

\begin{proof}
The displayed formula is $\lambda^2=2$ together with
$\lambda(\lambda-1)=2-\sqrt2$. For $RR$, survival of the first move requires
$x<r$ and of the second $\lambda x<r$, that is $x<1/2$, which is the stronger
condition. The case of $LL$ is the mirror image under $x\mapsto1-x$.
\end{proof}

Thus on even time scales the game contains the binary shift, and a block of type
$RR$ or $LL$ destroys only the single point $1/2$, in place of an interval of
length $g\approx0.293$.

\begin{proposition}\label{prop:closed}
Let $0<x<1$ and let $n\geq1$. Under blocks of type $RR$ and $LL$ the coordinate
$b$ doubles, and the position after $n$ such blocks equals
$\fp{2^{\,n-1}b\sqrt2}$, where $b$ is the coordinate at the start.
\end{proposition}

A knot in this regime is therefore described by a single integer, and the
sequence of blocks it requires is read off the binary expansion of
$\fp{b\sqrt2}$, a digit $0$ calling for $RR$ and a digit $1$ for $LL$.

\begin{definition}
For a configuration $S$ let $T(S)$ be the number of blocks of type $RR$ or $LL$
that $S$ survives entire, equivalently the length of the common binary prefix of
the positions of its knots.
\end{definition}

Since $\sqrt2$ has continued fraction $[1;2,2,2,\dots]$, its convergent
denominators are the Pell numbers and $\|q\sqrt2\|\asymp q^{-1}$. Taking
$b_i=iq$ for a Pell denominator $q$ places the values $\fp{b_i\sqrt2}$ within
$O(k/q)$ of one another, so that $k$ knots at those positions share about
$\log_2(q/k)$ blocks. Arbitrarily many knots can therefore share arbitrarily
long block words, and there is no positional obstruction to unbounded growth.

The obstruction lies in reaching such positions. A newly created knot has
$b=0$; its first move carries it to $b=\pm1$, and each subsequent block sends
$b\mapsto2b+2D$ with $D=\varepsilon-\delta\in\{0,\pm1\}$ determined by the side
of the deleted interval on which the knot falls. Steering a knot to a prescribed
$b$ costs about $\log_2|b|$ blocks, during which the same moves act on every
other knot. The following table, from exhaustive search to depth $30$, records
the largest $T$ attainable by a configuration of $k$ knots.

\begin{center}
\begin{tabular}{lccccc}
\toprule
$k$ & 3 & 4 & 5 & 6 & 7\\
\midrule
$\max T$ & 7 & 3 & 2 & 2 & 2\\
\bottomrule
\end{tabular}
\end{center}

Each additional knot roughly halves the shared prefix, which is the behaviour to
be expected of positions obeying the exponential separation guaranteed by
Lemma~\ref{lem:overlap}. Moreover $T$ affords no freedom of manoeuvre: within a
block of type $RR$ or $LL$ the choice between them is forced by the shared
digit. The property that prevents the collapse seen at $\varphi$ is thus also
the property that prevents clustered configurations from forming.

The record configurations found at $\lambda=\sqrt2$ have coordinates $b$, up to
sign, given by the first $5$, $7$ and $8$ terms of
\[
  0,\ 2,\ 6,\ 22,\ 54,\ 750,\ 13110,\ 26134
\]
for $5$, $7$ and $8$ knots respectively, and $0,2,6,22,54,408$ for $6$ knots.
The configurations attaining $5$, $7$ and $8$ knots are nested as sets of
points, each obtained from the previous by adjoining knots of much larger
$|b|$, and the configuration attaining $6$ knots is the reflection of the first
of these together with one further knot. From the configuration of eight knots,
however, no run of length at most $16$ produces a ninth or returns a superset,
and no prefix of length at most $20$ prepended to the corresponding run of
length $31$ produces a ninth.

\subsection{The binary square}

Because $\lambda^2=2$, the representation of Proposition~\ref{prop:itinerary}
splits by parity into two ordinary binary expansions:
$\lambda^{-2i}=2^{-i}$ and $\lambda^{-(2i-1)}=\sqrt2\,2^{-i}$. Set
\[
  \alpha_n=\sum_{i\geq1}\varepsilon_{n+2i}2^{-i}
  \;=\;0.\varepsilon_{n+2}\varepsilon_{n+4}\cdots,
  \qquad
  \beta_n=\sum_{i\geq1}\varepsilon_{n+2i-1}2^{-i}
  \;=\;0.\varepsilon_{n+1}\varepsilon_{n+3}\cdots,
\]
both in $[0,1]$.

\begin{proposition}\label{prop:square}
For every $n$,
\[
  x_n=(\sqrt2-1)\bigl(\alpha_n+\sqrt2\,\beta_n\bigr),
  \qquad
  (\alpha_{n+1},\beta_{n+1})=\bigl(\{2\beta_n\},\,\alpha_n\bigr),
  \qquad
  \varepsilon_{n+1}=\lfloor 2\beta_n\rfloor .
\]
\end{proposition}

\begin{proof}
The first is Proposition~\ref{prop:itinerary} with the parity split, since
$(\lambda-1)(\alpha_n+\sqrt2\beta_n)$ is the displayed tail. The second and
third are immediate from the definitions, the shift of the odd-indexed digits
being $\{2\beta_n\}$ and their leading digit $\lfloor2\beta_n\rfloor$.
\end{proof}

So the state of a knot is a point of the unit square, one move shifts and swaps
the coordinates, and two moves apply the binary shift to each: this is
Proposition~\ref{prop:blocks} in another form, the pairing of moves being the
splitting by parity. The position is the linear form
$x=(\sqrt2-1)(\alpha+\sqrt2\beta)$.

Unlike the three cells at $\lambda=3/2$, the deleted intervals here do not tile,
and $(0,1)$ falls into five regions, in order: destroyed by $L$; safe; destroyed
by $M$; safe; destroyed by $R$. For a knot in a given region with a given
required branch $\varepsilon$, the moves that both spare it and induce that
branch are as follows.

\begin{center}
\begin{tabular}{lcc}
\toprule
region & $\varepsilon=0$ & $\varepsilon=1$\\
\midrule
$(0,g)$ & $\{M,R\}$ & ---\\
$(g,r/2)$ & $\{M,R\}$ & $\{L\}$\\
$(r/2,1-r/2)$ & $\{R\}$ & $\{L\}$\\
$(1-r/2,r)$ & $\{R\}$ & $\{L,M\}$\\
$(r,1)$ & --- & $\{L,M\}$\\
\bottomrule
\end{tabular}
\end{center}

Every live knot therefore demands one of the four sets $\{L\}$, $\{R\}$,
$\{L,M\}$, $\{M,R\}$; all of them survive exactly when the demanded sets have a
common element, and a knot is created exactly when that element is $M$. Since
$\{L\}\cap\{M,R\}=\emptyset$ and $\{R\}\cap\{L,M\}=\emptyset$, the whole
blocking structure at this parameter is contained in that four-element list.

\begin{proposition}\label{prop:squaresurv}
Write $s_n=\alpha_n+\sqrt2\beta_n\in[0,1+\sqrt2]$. A knot admits some move
sparing it and inducing its required branch unless either $s_n<1/\sqrt2$ and
$\beta_n\geq1/2$, or $s_n>(2+\sqrt2)/2$ and $\beta_n<1/2$. A single knot's
itinerary is realisable by some run precisely when neither occurs at any time.
\end{proposition}

\begin{proof}
The two exceptional cases are the entries marked --- in the table, since
$x<g$ is equivalent to $s<1/\sqrt2$ and $x>r$ to $s>(2+\sqrt2)/2$, and
$\varepsilon=\lfloor2\beta\rfloor$.
\end{proof}

The contrast with Section~\ref{sec:threehalves} is now structural rather than
descriptive. At $\lambda=3/2$ the game is a one-dimensional map of multiplier
$3/2$ on an interval; at $\lambda=\sqrt2$ it is a two-dimensional binary shift
on a square with linear killing regions. This is the same distinction as the one
recorded in Section~\ref{sec:nonpisot}: $\Z[\sqrt2]$ is a lattice of rank two,
while $\Z[3/2]=\Z[1/2]$ is not a lattice at all.

\section{Questions}\label{sec:questions}

\begin{question}\label{q:kind}
Is there a $\lambda\in(1,2)$ and a periodic sequence $v$ with $v_p=M$ that is
kind? By Proposition~\ref{prop:kindyield} a single instance gives
$N_\lambda=\infty$. This is the sharpest form in which the unboundedness
question is currently posed: it asks for one finite word.
\end{question}

\begin{question}
Is $N_{\sqrt2}$ bounded? Is $N_{3/2}$? By the remark following
Corollary~\ref{cor:noperiodic}, at these two parameters $1/2$ is not a periodic
point for any itinerary, so a kind periodic sequence would settle
Question~\ref{q:kind} affirmatively at a parameter of interest.
\end{question}

\begin{question}
Is the Pisot condition necessary as well as sufficient for boundedness? Is
there a non-Pisot $\lambda$, or a $\lambda$ that is not an algebraic integer,
for which $N_\lambda$ is bounded?
\end{question}

\begin{question}
For $\lambda=\sqrt2$, does $\max\{T(S):|S|=k\}$ remain at least $1$ for every
$k$? A negative answer would say that beyond some count the knots cannot agree
on which side of $1/2$ they lie, and would be evidence for boundedness.
\end{question}

\begin{question}
What happens for Salem $\lambda$, whose conjugates lie on the unit circle? The
orbit of $1/2$ should then be infinite but confined, a regime intermediate
between those of Sections~\ref{sec:pisot} and \ref{sec:nonpisot}.
\end{question}

\begin{question}
Is $N_\lambda$ bounded for all $\lambda$ in some interval $(\lambda^*,2)$? The
newborn knot is carried to within $(2-\lambda)/2$ of an endpoint, after which
$\min(x,1-x)$ is multiplied by $\lambda$ at each move, so about
$\log(1/(2-\lambda))/\log\lambda$ moves are needed to return to the middle of
the strip. This suggests $N_\lambda\asymp\log\bigl(1/(2-\lambda)\bigr)$ as
$\lambda\uparrow2$. Such a result would be the first boundedness statement not
obtained from Theorem~\ref{thm:pisot}, and would have to be metric rather
than arithmetic.
\end{question}

\begin{question}\label{q:mahlerform}
Is there a periodic $c\in\{0,1,2\}^{\mathbb N}$ containing a $1$, and an index
$t$ with $c_t=1$, such that the orbit of $w=1$ under
$w\mapsto\tfrac32w-[\,\lfloor\tfrac32w\rfloor>c_s\,]$ started at time $t$
satisfies $\lfloor\tfrac32w_s\rfloor\neq c_{s+1}$ for every $s\geq t$? An
affirmative answer gives $N_{3/2}=\infty$, by
Proposition~\ref{prop:immortal32}. This is Question~\ref{q:kind} at
$\lambda=3/2$, stated without reference to the game.
\end{question}

\begin{question}
At $\lambda=(1+\sqrt2)/2$ the record lengths satisfy
$d(k)=\max(2k-1,\lfloor k^2/4\rfloor+1)$ at every $k$ where they are known
exactly. Does the formula persist? More broadly, is any superlinear lower bound
on $d_\lambda(k)$ provable at any parameter? Such a bound would have to
quantify the obstruction to creating knots as the count grows, and none is
known.
\end{question}

\begin{question}\label{q:ae}
Prove that the kind endpoints are dense in $(0,1)$ for a single non-Pisot
$\lambda$ --- by Theorem~\ref{thm:density0} this alone gives
$N_\lambda=\infty$ --- or prove Hypothesis \textup{(D$_\lambda$)} of
Section~\ref{sec:density} for almost every $\lambda\in(1,2)$. By
Theorem~\ref{thm:density} either would give unbounded knot counts, and the
almost-everywhere version would establish that $\{\lambda:N_\lambda<\infty\}$
has measure zero, as the conjecture that this set is exactly the Pisot numbers
predicts. Section~\ref{sec:transversality} supplies the certified first steps of a
parameter-exclusion approach on a window containing $3/2$, including an
unconditional exponential lower bound on the number of kind words; what
remains is the half that abundance does not touch --- that the endpoints of
those words spread out rather than pile up --- and the exclusion iteration
assembling the estimates.
\end{question}

\begin{question}
The least lengths for $\lambda=3/2$ begin $1,3,5,9,19,23,52$, with
$d_{3/2}(8)\geq57$. The growth is irregular: the run of increments
$2,2,4,10,4,29$ fits neither a clean exponential (an earlier draft of this
paper suggested $2^k$, which $d(7)=52<128$ refutes) nor the quadratic law that
fits all known values at $(1+\sqrt2)/2$. Determine the order of growth of
$d_{3/2}(k)$.
\end{question}

\appendix

\section{Verification}\label{app:verify}

Every computation reported above was carried out in exact arithmetic. Positions
lie in $\tfrac12\Z[\lambda]$ and were represented by integer coordinate vectors;
comparisons against the endpoints of the deleted intervals reduce to the sign of
an element of $\Z[\lambda]$, evaluated exactly for quadratic $\lambda$ and, for
the plastic number, by an exact test for vanishing together with a
high-precision evaluation of the sign, which suffices because the relevant set
of points is finite.

Two independent implementations were compared for each of the three quadratic
parameters: one in the coordinates just described, and one simulating the strip
literally as a moving subinterval of the line with absolute knot positions and
no rescaling. The two agree on all $88{,}572$ runs of length at most $10$ for
$\lambda=3/2$, all $29{,}523$ runs of length at most $9$ for $\lambda=\sqrt2$,
and all $9{,}840$ runs of length at most $8$ for $\lambda=\varphi$. The circle
normal form of Proposition~\ref{prop:circle} was checked against the interval
formulation on all $265{,}719$ runs of length at most $11$ for $\lambda=3/2$,
with the seam remaining at $0$ throughout.

Lemma~\ref{lem:distinct}, Lemma~\ref{lem:suffix}, Theorem~\ref{thm:mono},
Proposition~\ref{prop:big}, Theorem~\ref{thm:pisot}, Theorem~\ref{thm:phi} with
the table of Appendix~\ref{app:phi}, Proposition~\ref{prop:littlewood},
Corollary~\ref{cor:noperiodic}, Proposition~\ref{prop:kindyield},
Proposition~\ref{prop:blocks} and Proposition~\ref{prop:closed} have been
formally verified in Lean~4 against Mathlib, with a semantic audit confirming
that every public theorem depends only on \texttt{propext},
\texttt{Classical.choice} and \texttt{Quot.sound}. For the plastic number only
the hypotheses of Proposition~\ref{prop:plastic} are certified, namely that
$\rho$ is Pisot and that its orbit is therefore finite; the three numerical
assertions are not, the kernel-checked closure computation at that size being
infeasible without compiled evaluation. Everything listed below as out of scope
remains a computation only.

A second round of verification covers Section~\ref{sec:gaps}.
Lemmas~\ref{lem:order}, \ref{lem:gaplaw} and~\ref{lem:service} are certified as
stated; Proposition~\ref{prop:lowerbound} is certified in the form the
scheduling argument consumes (no birth is followed immediately by $M$, and a
word of length $n$ carries at most $(n+1)/2$ surviving births);
Theorem~\ref{thm:schedule} is certified with the logarithm-free hypothesis
$1\leq\lambda^W\delta$ in place of $W=\lceil\log_\lambda(1/\delta)\rceil$ and
with $\lceil W/2\rceil$ as natural division; Corollary~\ref{cor:effective} is
certified at the golden ratio, $N_\varphi\leq9$, with no numerics, via the
exact identity $\varphi^5(\varphi-\tfrac32)=\varphi^2/2$; and, resolving part
of the plastic gap above, $N_\rho\leq34$ is certified: the $153$-point orbit,
its closure under both maps, the membership of $1/2$, and the gap bound
$\delta_0\geq239/100000$ are all checked by kernel reduction over integer
coordinate triples, with $\rho$ bracketed between rationals scaled to integer
inequalities. The configuration closure and the exact maximum $7$ remain
computations only. The formal proof of Theorem~\ref{thm:schedule} differs from
the one printed here and is simpler: an induction on the window word that, at
each straddling move, deletes the largest knot below the deleted interval --
no ages, no injection.

The third round also certifies the transversality material of
Section~\ref{sec:transversality}: Proposition~\ref{prop:restricted} by a
$27$-cell certificate closed by kernel reduction, with all arithmetic over the
common denominator $1024000000$ and the cell decomposition proved to tile the
window, and Proposition~\ref{prop:paircount} with Mathlib's Lebesgue measure,
instantiated at the certified window. Later rounds certify Lemma~\ref{lem:survivorset} (with the cells as an
explicit list carrying a disjointness invariant), Lemma~\ref{lem:permanent},
Lemma~\ref{lem:oneperstep}, Theorem~\ref{thm:compactness} -- where the
compactness step is carried out with a non-principal ultrafilter rather than a
diagonal subsequence, and $N_\lambda=\infty$ is rendered as unboundedness of
$N_\lambda(n)$ -- Theorem~\ref{thm:density0} with its hypothesis explicit in
every statement, Corollary~\ref{cor:deficit}, Lemma~\ref{lem:free},
Remark~\ref{rem:candidates}'s cell identity, and
Proposition~\ref{prop:expcount}. The counterexample recorded in
Lemma~\ref{lem:free} for $\lambda\geq3/2$ was found in the course of that
verification, and corrects an earlier version of this paper which asserted the
converse without restriction. For Proposition~\ref{prop:twostep} the runs and
the containments of Figure~\ref{fig:records32} are certified for the words
exhibited; the exhaustiveness claim, which needed millions of retained states
at $k=7$, is not, and remains a computation. Also certified, beyond what earlier drafts of this appendix credited: the
Littlewood identity and Corollary~\ref{cor:noperiodic}; the block
reformulation of Section~\ref{sec:sqrt2} in the form
$x\mapsto\{2^n(a+b\sqrt2)/2\}$ along blocks of two moves; and, at the plastic
number, the full configuration closure of $25{,}525$ states with
$N_\rho\leq7$ and the values $d_\rho(4)=7$, $d_\rho(7)=28$ --- the computation
reported infeasible after round 1, later carried through. A further round certifies the whole of Section~\ref{sec:threehalves}:
Propositions~\ref{prop:ternary}, \ref{prop:base32} and~\ref{prop:mahler}, the
equivalence of Proposition~\ref{prop:translation} in both directions, and the
fact that reachable positions at $3/2$ avoid the ternary cell boundaries. It
also certifies Proposition~\ref{prop:kindyield}, the cylinder counts of
Proposition~\ref{prop:kinddim} (the dimension statements are deliberately
\emph{not} formalised, and the development says so, since the count alone does
not yield them), the quantitative density criterion with the explicit bound
recorded in Section~\ref{sec:density}, and the sharpness witnesses of
Section~\ref{sec:transversality}, found by untrusted search programs and
re-verified by the kernel. One census finding: the itinerary identity of
Proposition~\ref{prop:itinerary} had never been formalised, and is now proved
for every $\lambda>1$. The audit after this round reports $1198$ theorems,
axioms confined to the same three.

The audit at the end of the preceding rounds reported $1084$ theorems.

A fourth round certifies
Lemmas~\ref{lem:nojump} and~\ref{lem:goodchild} (with each crossing inequality
proved equivalent to $\lambda^2\geq\lambda+1$),
Proposition~\ref{prop:continuum} -- the injection valid for all
$\lambda\in(1,\varphi)$, since only finiteness of the forced return time is
required -- and Proposition~\ref{prop:commonwindow}; after it the audit
reports $636$ theorems, axioms confined to the same three. One plan recorded in earlier commissions --- a probabilistic bound on the
return time of the discarded child, feeding a renewal inequality --- was
deliberately fenced off from formalisation because it was never proved on
paper, and it remains absent; the exponential count of
Proposition~\ref{prop:expcount} was obtained instead by the covering argument
described with it, which needs neither ingredient.

A third round likewise certifies the two cubic parameters of
Section~\ref{sec:pisot}: at the tribonacci constant, the seven-point orbit and
its closure, the $20$ reachable configurations, the maximum $3$ and the values
$d=1,3,7$; at the supergolden ratio, the $43$-point orbit, the maximum $4$ and
$d=1,3,5,11$, where what is certified of the $412$-entry configuration list is
what the argument consumes -- closure under the three moves and the bound
$4$ -- rather than its irredundancy. The audit standard is unchanged
throughout: $524$ theorems after the third round, axioms confined to the same
three, no \texttt{sorry}, \texttt{admit} or \texttt{native\_decide}.

Proposition~\ref{prop:itinerary} was checked at $\lambda=3/2$, $\sqrt2$ and
$\varphi$ against trajectories of length $200$, and
Propositions~\ref{prop:square} and \ref{prop:squaresurv} on $100$ steps of an
exact trajectory at $\lambda=\sqrt2$, in each case to within the truncation
bound and without a discrepancy.

The exact values $d(10)=26$ and $d(11)=31$ at $\lambda=(1+\sqrt2)/2$ were
established by a targeted form of the exhaustive search, asking whether $k$
knots are attainable by depth $H$ and pruning with three reductions, each of
which preserves the answer: identification of mirror-image configurations under
$x\mapsto1-x$, which exchanges $L$ and $R$; subset domination as above; and the
birth-spacing bound underlying Proposition~\ref{prop:lowerbound}, by which a
configuration of $j$ knots at depth $n$ reaches $k$ knots by depth $H$ only if
$j\geq k-\lceil(H-n)/2\rceil$. On the known case, nine knots by depth $20$, the
pruned search retains ten configurations at depth $18$ where the unpruned
search retains $1{,}192{,}182$ at depth $20$, and returns the same answer. The
searches for ten knots by depth $25$ and eleven by depth $30$ then run to
completion, and beam searches supply witnesses at depths $26$ and $31$. The
corresponding search for twelve knots by depth $36$ exhausts memory near depth
$22$, where its threshold is still weak.

The value $d_{3/2}(7)=52$ was fixed by the same targeted method run at
$\lambda=3/2$: complete searches exclude seven knots at every depth from $43$
to $51$, and a witness of length $52$ exists. For the horizon-$43$ search the
retained layers peak at $106{,}569$ configurations, against layers in the
millions for the unthresholded search. Complete searches likewise exclude
eight knots at depths $54$, $55$ and $56$, which with
Corollary~\ref{cor:recursive} gives $d_{3/2}(8)\geq57$; at horizon $57$ and
beyond the retained layers exceed three million states near depth $46$ and the
method fails.


The lemmas of Section~\ref{sec:gaps} were tested exactly, in
$\Z[\sqrt2]$, on a $63$-move run at $\lambda=(1+\sqrt2)/2$ carrying fourteen
simultaneous knots: the gap law held for all $1{,}091$ surviving pairs
($907$ same-side, $184$ straddling) with no violations, and the order of the
knots was preserved at every step. The identity for the deficit
$1-(\mathrm{spread})$ that follows from the gap law held at $58$ of the $60$
steps at which two or more knots survived, the two exceptions being exactly
the steps at which a knot died, where no identity is asserted.

The tribonacci and supergolden computations of Section~\ref{sec:pisot} were
run in $\Z[\lambda]$ with doubled integer numerators $(A,B,C)$ for
$x=(A+B\lambda+C\lambda^2)/2$; signs of nonzero elements were decided by
$60$-digit evaluation guarded by the threshold $10^{-30}$, with exact zeros
detected as vanishing coefficient triples, and the guard never tripped. The
coverage depths $m^*(L)$ of Section~\ref{sec:density} at $\lambda=3/2$ were
computed in exact integer arithmetic. The eighteen-parameter survey of that
section is the one floating-point computation cited in this paper: positions
were deduplicated at $10^{-11}$, which inflated the plastic orbit to $155$
points against the exact $153$; its conclusions are stated only qualitatively.

The reformulations of Section~\ref{sec:threehalves} were checked against the
original rules in exact arithmetic: Proposition~\ref{prop:ternary} on all
$797{,}160$ runs of length at most $12$, and Propositions~\ref{prop:base32} and
\ref{prop:mahler} on $300$ trajectories of length $60$ and on $18{,}000$
individual steps respectively, in each case without a discrepancy.

The exhaustive depths quoted are those independently re-established with
subset-domination pruning: depth $42$ for $\lambda=3/2$ and depth $38$ for
$\lambda=\sqrt2$, at which the largest counts are six and eight respectively.

Searches described as exhaustive enumerate all reachable configurations at each
depth, with configurations that are subsets of others discarded where noted;
such discarding is justified because survival is a condition on individual
points, so that $A\supseteq B$ implies the same containment after any move.
Searches described as beam searches retain a bounded number of configurations at
each depth and yield lower bounds on $N_\lambda$ only.

The search of Section~\ref{sec:kind} illustrates why exact arithmetic is not a
formality here. Carried out instead by evaluating the orbit to fifty decimal
places, it reports several hundred realisable itineraries admitting a periodic
run with an $M$. Every one of them is spurious. The quantity being tested at the
final step is $x_{q-1}-r/2$ or $x_{q-1}-(1-r/2)$, which by
Lemma~\ref{lem:distinct} is exactly zero, so a floating evaluation returns whichever
sign the rounding supplies. The same effect makes a direct simulation of the
periodic runs at algebraic $\lambda$ overcount badly: at $\varphi$ it reports
twenty-seven knots and at the plastic number twenty, where
Theorem~\ref{thm:phi} and Proposition~\ref{prop:plastic} give two and seven. The
coincidences responsible for the small exact answers are precisely what
floating arithmetic destroys.

\section{Exploratory computations}\label{app:explore}

The computations of this appendix guided the investigation without entering
any proof. They are floating-point unless stated otherwise, and are recorded
so that the reader can see which quantities were suggestive and which were
merely consistent.

\subsection{A first moment}
Suppose the three moves are chosen independently and uniformly. A move deletes
a closed interval of length $g$ wherever it is placed, so a knot at a position
independent of the move survives with probability $r$, whichever move is
played. Births occur at rate $1/3$, so the expected number of knots present
after many moves is
\[
  \sum_{a\geq0}\tfrac13r^a=\frac{1}{3(1-r)}=\frac{\lambda}{3(\lambda-1)} ,
\]
which is $1$ at $\lambda=3/2$ and $1.94$ at $\lambda=(1+\sqrt2)/2$. This
baseline is a useful sanity check on simulations, and it makes plain how far
the maxima exceed the typical behaviour: the game is entirely about the tail
of a distribution whose mean is a small constant. It also shows that no
counting argument of this kind can distinguish Pisot from non-Pisot
parameters, since the mean does not see the arithmetic at all.

\subsection{Backward counts}
Reading a word backwards and applying the two inverse branches to $1/2$
produces, at every non-Pisot parameter examined, $2^n$ distinct points after
$n$ steps: backwards there are no coincidences. At $\varphi$ the count is
$F(n+3)-1$ for the Fibonacci numbers $F$, growing like $\varphi^n$ rather than
$2^n$ --- the same collapse recorded in Figure~\ref{fig:reachable} for the
forward direction, in its backward form.

\subsection{The transform along a survey}
Writing $|\widehat{\nu_r}(\xi)|$ for the cosine product of
Section~\ref{sec:conjugate} and taking the maximum of its logarithm over
$\xi\in[140,210]$ --- an upper envelope, computed on a fine grid so that
narrow peaks are not missed --- gives
\[
\varphi:-1.37,\quad 3/2:-2.44,\quad \rho:-2.91,\quad
\sqrt2:-4.71,\quad (1+\sqrt2)/2:-7.34 .
\]
The ordering is instructive and slightly treacherous. Half the silver ratio
decays fastest, consistent with its being the most heavily overlapped of the
five; but $3/2$, which is not Pisot, sits \emph{above} the plastic number,
which is. At moderate frequencies a slowly decaying transform can exceed a
Pisot floor, so the value of $|\widehat{\nu_r}|$ is not by itself diagnostic;
only its behaviour as $\xi\to\infty$ is. The same caution applies in real
space, where the plastic density looks smooth at coarse resolution
(Figure~\ref{fig:densities}).

\subsection{Beam searches}
Upper bounds on $d_\lambda(k)$ marked as such in Section~\ref{sec:nonpisot}
come from beam searches: breadth-first exploration retaining a fixed number of
configurations per level, ranked by knot count and then by spread. These
establish upper bounds only, since every configuration retained is genuinely
reachable; the lower bounds quoted alongside them come from exhaustive search
or from Corollary~\ref{cor:recursive}. Beam widths between $10^4$ and
$5.5\times10^4$ were used, and the searches that stall do so through memory
rather than time.

\section{The golden ratio automaton}\label{app:phi}

The twelve reachable configurations of Theorem~\ref{thm:phi} and their
transitions are as follows, with $\emptyset\mapsto\{p_3\}$ under $M$ and
$\emptyset\mapsto\emptyset$ under $L$ and $R$.

\begin{center}
\begin{tabular}{lccc}
\toprule
 & $L$ & $M$ & $R$\\
\midrule
$\{p_1\}$ & $\emptyset$ & $\{p_2,p_3\}$ & $\{p_2\}$\\
$\{p_2\}$ & $\emptyset$ & $\{p_3\}$ & $\{p_3\}$\\
$\{p_3\}$ & $\{p_1\}$ & $\{p_3\}$ & $\{p_5\}$\\
$\{p_4\}$ & $\{p_3\}$ & $\{p_3\}$ & $\emptyset$\\
$\{p_5\}$ & $\{p_4\}$ & $\{p_3,p_4\}$ & $\emptyset$\\
$\{p_1,p_3\}$ & $\{p_1\}$ & $\{p_2,p_3\}$ & $\{p_2,p_5\}$\\
$\{p_1,p_4\}$ & $\{p_3\}$ & $\{p_2,p_3\}$ & $\{p_2\}$\\
$\{p_2,p_3\}$ & $\{p_1\}$ & $\{p_3\}$ & $\{p_3,p_5\}$\\
$\{p_2,p_5\}$ & $\{p_4\}$ & $\{p_3,p_4\}$ & $\{p_3\}$\\
$\{p_3,p_4\}$ & $\{p_1,p_3\}$ & $\{p_3\}$ & $\{p_5\}$\\
$\{p_3,p_5\}$ & $\{p_1,p_4\}$ & $\{p_3,p_4\}$ & $\{p_5\}$\\
\bottomrule
\end{tabular}
\end{center}

\begin{thebibliography}{9}

\bibitem{Bertin1992}
M.-J. Bertin, A. Decomps-Guilloux, M. Grandet-Hugot, M. Pathiaux-Delefosse and
J.-P. Schreiber,
\emph{Pisot and Salem Numbers},
Birkh\"auser, Basel, 1992.

\bibitem{Erdos1939}
P. Erd\H{o}s,
\emph{On a family of symmetric Bernoulli convolutions},
Amer. J. Math. \textbf{61} (1939), 974--976.

\bibitem{ErdosJooKomornik1990}
P.~Erd\H{o}s, I.~Jo\'o and V.~Komornik,
\newblock Characterization of the unique expansions
$1=\sum q^{-n_i}$ and related problems,
\newblock \emph{Bull.\ Soc.\ Math.\ France} \textbf{118} (1990), 377--390.

\bibitem{Garsia1962}
A. M. Garsia,
\emph{Arithmetic properties of Bernoulli convolutions},
Trans. Amer. Math. Soc. \textbf{102} (1962), 409--432.

\bibitem{Hochman2014}
M. Hochman,
\emph{On self-similar sets with overlaps and inverse theorems for entropy},
Ann. of Math. (2) \textbf{180} (2014), 773--822.

\bibitem{Koksma1935}
J.~F. Koksma,
\newblock Ein mengentheoretischer Satz \"uber die Gleichverteilung modulo Eins,
\newblock \emph{Compositio Math.} \textbf{2} (1935), 250--258.

\bibitem{Mahler1968}
K. Mahler,
\emph{An unsolved problem on the powers of $3/2$},
J. Austral. Math. Soc. \textbf{8} (1968), 313--321.

\bibitem{Solomyak1995}
B.~Solomyak,
\newblock On the random series $\sum\pm\lambda^n$ (an Erd\H{o}s problem),
\newblock \emph{Ann.\ of Math.\ (2)} \textbf{142} (1995), 611--625.

\bibitem{PeresSolomyak1996}
Y.~Peres and B.~Solomyak,
\newblock Absolute continuity of Bernoulli convolutions, a simple proof,
\newblock \emph{Math.\ Res.\ Lett.} \textbf{3} (1996), 231--239.

\bibitem{Salem1944}
R. Salem,
\emph{A remarkable class of algebraic integers. Proof of a conjecture of
Vijayaraghavan},
Duke Math. J. \textbf{11} (1944), 103--108.

\end{thebibliography}

\end{document}
